\documentclass[11pt]{amsart}
\usepackage{lmodern}
\usepackage{amsmath,amssymb,amsthm}
\usepackage[margin=1.05in]{geometry}
\usepackage{booktabs}
\usepackage[hidelinks]{hyperref}
\setlength{\emergencystretch}{2em}
\renewcommand{\labelitemi}{\(\bullet\)}

\newtheorem{theorem}{Theorem}[section]
\newtheorem{lemma}[theorem]{Lemma}
\newtheorem{proposition}[theorem]{Proposition}
\newtheorem{corollary}[theorem]{Corollary}
\newtheorem{conjecture}[theorem]{Conjecture}
\theoremstyle{definition}
\newtheorem{definition}[theorem]{Definition}
\newtheorem{problem}[theorem]{Problem}
\theoremstyle{remark}
\newtheorem{remark}[theorem]{Remark}

\newcommand{\Z}{\mathbb{Z}}
\newcommand{\Q}{\mathbb{Q}}
\newcommand{\Zp}{\mathbb{Z}_p}
\newcommand{\vp}{v_p}
\newcommand{\bin}[2]{\binom{#1}{#2}}

\title[Sharp Ap\'ery defect laws]{Sharp Ap\'ery defect laws,\\ optimal tower exponents, and infinite counterexamples}
\date{November 14, 2026 --- revised proof draft}

\begin{document}

\begin{abstract}
Let $A(n)$ and $B(n)$ be the Ap\'ery sequences attached to $\zeta(3)$ and
$\zeta(2)$. We transfer their central-binomial rigidity to three further
questions. For
\[
u(n)=\sum_{k=0}^{n-1}\binom nk^2\binom{n+k-1}k^2
\quad\text{(OEIS A361713)},
\]
we prove the all-level bound
\[
v_p\bigl(u(p^r)-u(p^{r-1})\bigr)\ge\min(5r,3r+5)
\qquad(p\ge7,\ r\ge1).
\]
This exponent is best possible at every level: the fixed prime $11$
attains equality for every $r$. A first-defect formula explains the
change of exponent and supplies an infinite family contradicting the
older, already refuted $p^{4r+1}$ proposal.
Independently, the smooth-index conjecture on A357567 is false:
if $F(n)=3^{42}A(n)^{25}-5^{25}B(n)^{42}$, then
$v_5(F(3\cdot5^r))=25$ for every $r\ge0$, so all $r\ge6$ give
counterexamples to divisibility by $(3\cdot5^r)^5$.
Finally, an adjacent telescoping identity proves Bala's second conjecture
on A361717 for every prime other than $3$, with a fifth-order refinement.
The preceding rigidity proofs are included to make the dependencies
explicit in one source file. No complete proof of Swisher's stronger
$(C.3)$ congruence is claimed; its nonzero boundary obstruction is recorded.
\end{abstract}

\maketitle

\section*{Guide to the transfer results}
The main transfer theorem is Theorem~\ref{thm:sharp-J4}. It replaces the
false exponent $4r+1$ for A361713 by the optimal universal exponent
$\min(5r,3r+5)$ and determines the first nonzero defect at $p=11$.
The existing counterexample at $p=7,r=6$ is credited to Adamczewski
and is recorded on the OEIS page \cite{A361713,Adam}; the contribution
here is the all-level replacement and exact sharpness, not the first
disproof of that conjecture.

Theorem~\ref{thm:transfer-counterexample-five} gives a separate infinite
counterexample family to A357567 conjecture~4 \cite{A357567}. Its proof
is elementary and independent of the enhanced rigidity theorem.
The final adjacent-identity section resolves A361717 conjecture~2
\cite{A361717}. These statements have different logical statuses:
a corrected optimal theorem, a disproof, and a direct proof, respectively.

The foundational material follows first, including the complete proofs
of all rigidity inputs used in the transfer. A finite verification
package accompanies this research proof draft. Neither the computations
nor the present source review constitute external peer review or a
certificate of historical priority.

\section{Rigidity foundations}

\subsection{Bala's conjectures}
Throughout, $A(n)$ denotes the Ap\'ery numbers for $\zeta(3)$ (OEIS A005259) and $B(n)$ the
Ap\'ery numbers for $\zeta(2)$ (OEIS A005258):
\[
A(n)=\sum_{k=0}^n\bin nk^2\bin{n+k}k^2,\qquad B(n)=\sum_{k=0}^n\bin nk^2\bin{n+k}k .
\]
Classical supercongruences (Gessel, Beukers, Coster \cite{Ge,Be,Co}) give
$u(mp^r)\equiv u(mp^{r-1})\pmod{p^{3r}}$ for $u\in\{A,B\}$, $p\ge5$; in particular
$A(p)\equiv5$, $A(p-1)\equiv1$, $B(p)\equiv3$, $B(p-1)\equiv1 \pmod{p^3}$.
In 2021--2022 P.~Bala recorded in the OEIS a family of conjectures asserting that certain
combinations of these numbers satisfy congruences two or three powers of $p$ stronger
\cite{OEIS}. With
\[
L_r=\begin{cases}5,& r=1,\\ 3r+3,& r\ge 2,\end{cases}
\]
they read as follows; here $p\ge5$ is prime, and statement (P) is the prime level $r=1$
while (T) is the tower $r\ge2$, in the form $u(p^r)\equiv u(p^{r-1})\pmod{p^{L_r}}$.
\begin{itemize}
\item A212334: $a(n)=\sum_{k=0}^{n-1}\bin nk\bin{n-1}k\bin{n+k-1}k^2$ $(n\ge1)$, which equals
$\tfrac1{12}(A(n)+7A(n-1))$. Bala conjectured (P) $a(p)\equiv 1\pmod{p^5}$ (Dec.\ 2021) and
(T) (Oct.\ 2022).
\item A352655: $\tfrac12(B(n)+B(n-1))$; (P) $\equiv2\pmod{p^5}$ and (T).
\item A357506: $B(n)^3B(n-1)$; (P) $\equiv 27\pmod{p^5}$ for all $p\ge3$, and (T).
\item A357567: $5A(n)-14B(n)$; (P), (T), and the same for the companion
$3^{42}A(n)^{25}-5^{25}B(n)^{42}$.
\item A357956: $5A(n-1)-2B(n-1)$; (P) $\equiv 3\pmod{p^5}$ and (T) at $n=p^r$.
\item A357959: $5A(n-1)+2B(n)$; (P), (T), and the same for the companion $A(n-1)^5B(n)^6$.
\end{itemize}
A further conjecture of Bala concerns A361717,
$\nu(n)=\sum_{k=0}^{n-1}\bin{n-1}k^2\bin{n+k}k$: Conjecture~1 there is
$\nu(p)\equiv0\pmod{p^4}$ for $p\ge5$.

\subsection{Status before this work}
The supplied November 11 draft reported that the Formal Conjectures repository
\cite{FC} formalises both A212334 conjectures with the tag \emph{research open}.
A public AI-assisted campaign
of 2026 \cite{Cam} did two things. First, it proved the prime-level A212334 congruence (P),
using the same product identity as our Lemma~\ref{lem:reps}. Second, it showed that every claim
above except those on A212334 and A361717 follows from three linear ``defect''
relations $R_{1,r},R_{2,r},R_{3,r}$ (recalled in \eqref{eq:R}). Those relations, at every level
$r\ge 1$, and the A212334 tower, were listed there as open in that draft's source review.
The present source check likewise found the displayed $B$-tower conjecture in the OEIS.
This limited search cannot prove novelty, and
Section~\ref{sec:status} states precisely what we claim.

\subsection{A rigidity problem}
For a prime $p$ and $r\ge1$ we write $N=p^r$, $M=p^{r-1}$ and define the central-binomial
defect $\Delta_r$ and the four Ap\'ery defects:
\[
\Delta_r=\bin{2N}{N}-\bin{2M}{M},\qquad
\begin{aligned}
\gamma_r&=A(N)-A(M), & \delta_r&=A(N-1)-A(M-1),\\
\alpha_r&=B(N)-B(M), & \beta_r&=B(N-1)-B(M-1).
\end{aligned}
\]
In this notation Bala's linear relations are
\begin{equation}\label{eq:R}
R_{1,r}:\ \alpha_r+\beta_r\equiv0,\qquad
R_{2,r}:\ 5\gamma_r-14\alpha_r\equiv0,\qquad
R_{3,r}:\ 5\delta_r-2\beta_r\equiv0\pmod{p^{L_r}},
\end{equation}
and A212334 (T) is $\gamma_r+7\delta_r\equiv 0\pmod{p^{L_r}}$.
All four defects and $\Delta_r$ are $\equiv0\pmod{p^{3r}}$ for $p\ge5$. We propose the
following question, which to our knowledge has not been asked before.

\begin{problem}[central-binomial rigidity]\label{prob}
For an Ap\'ery-type sequence $u$, do there exist rational numbers $c_u$ such that
$u(p^r)-u(p^{r-1})\equiv c_u\,\Delta_r\pmod{p^{L_r}}$ for all large primes $p$ and all $r\ge1$?
\end{problem}

A positive answer with constants $c_A,c_{A^-},c_B,c_{B^-}$ for the four defects would give all
of Bala's linear claims as linear algebra in these four rationals. The constants turn out to be
\[
(c_A,\ c_{A^-},\ c_B,\ c_{B^-})=\left(\tfrac72,\ -\tfrac12,\ \tfrac54,\ -\tfrac54\right),
\]
where $c_{A^-}$ refers to $A(n-1)$ and $c_{B^-}$ to $B(n-1)$.

\subsection{Results}
\begin{theorem}[Bala's A212334 conjectures]\label{thm:A212334}
Let $p\ge5$ be prime and $r\ge1$. Then $a(p^r)\equiv a(p^{r-1})\pmod{p^{L_r}}$.
That is, $a(p)\equiv1\pmod{p^5}$, and $a(p^r)\equiv a(p^{r-1})\pmod{p^{3r+3}}$ for $r\ge2$.
\end{theorem}

\begin{theorem}[rigidity for $\zeta(3)$]\label{thm:W}
Let $p\ge 7$ and $r\ge1$, or $p=5$ and $r\ge2$. Then
\[
\gamma_r\equiv\tfrac72\Delta_r,\qquad \delta_r\equiv-\tfrac12\Delta_r\pmod{p^{L_r}}.
\]
\end{theorem}

\begin{theorem}[rigidity for $\zeta(2)$ at prime level]\label{thm:Z1}
Let $p\ge7$. Then $\alpha_1\equiv\tfrac54\Delta_1$ and $\beta_1\equiv-\tfrac54\Delta_1\pmod{p^5}$.
\end{theorem}

\begin{corollary}[Bala's prime-level conjectures]\label{cor:prime}
For every prime $p\ge5$:
\begin{align*}
&\tfrac12(B(p)+B(p-1))\equiv2, && 5A(p)-14B(p)\equiv-17, && 5A(p-1)-2B(p-1)\equiv3,\\
&5A(p-1)+2B(p)\equiv 11, && 3^{42}A(p)^{25}-5^{25}B(p)^{42}\equiv0, && A(p-1)^5B(p)^6\equiv 3^6,
\end{align*}
all modulo $p^5$. Moreover $B(p)^3B(p-1)\equiv 27\pmod{p^5}$ for every prime $p\ge3$.
These are the prime-level statements on A352655, A357567 (1) and (3),
A357956 (1), A357959 (1) and (3), and A357506.
\end{corollary}

\begin{theorem}[Bala's Conjecture~1 on A361717]\label{thm:nu}
For every prime $p\ge 5$,
\[
\nu(p)=\sum_{k=0}^{p-1}\bin{p-1}k^2\bin{p+k}k\equiv0\pmod{p^4}.
\]
\end{theorem}

The prime-power towers follow from the following reduction and its completion below.

\begin{theorem}[harmonic comparison and tower reduction]\label{thm:towers}
Let $p\ge7$, $r\ge1$, or $p=5$, $r\ge2$. Put $N=p^r$, $M=p^{r-1}$, and let
$I_j(n)=\sum_{k=1}^{n-1}(n/k)^jO_k(n)E_k(n)$ be the odd sums of \eqref{eq:IJ}.
\begin{enumerate}
\item[(a)] $N H_{N-1}-M H_{M-1}\equiv \tfrac14\Delta_r\pmod{p^{L_r}}$.
\item[(b)] $I_3(N)\equiv I_3(M)\pmod{p^{L_r}}$.
\item[(c)] Put $\mathcal D_r=I_2(N)-I_2(M)-\tfrac14\Delta_r$. Then
$\alpha_r\equiv\tfrac54\Delta_r+\mathcal D_r$ and $\beta_r\equiv-\tfrac54\Delta_r-3\mathcal D_r
\pmod{p^{L_r}}$. Consequently $R_{1,r}\Leftrightarrow 2\mathcal D_r\equiv0$,
$R_{2,r}\Leftrightarrow 14\mathcal D_r\equiv 0$ and $R_{3,r}\Leftrightarrow 6\mathcal D_r\equiv0$,
modulo $p^{L_r}$.
\end{enumerate}
\end{theorem}

For $p\ge7$ and $r=1$, Theorem~\ref{thm:Z1} gives $\mathcal D_1\equiv0\pmod{p^5}$.
\begin{theorem}[completion of central-binomial rigidity]\label{conj:D}
For all primes $p\ge5$ and all $r\ge2$,
\[
\mathcal D_r\equiv0,\qquad
\alpha_r\equiv\tfrac54\Delta_r,\qquad
\beta_r\equiv-\tfrac54\Delta_r\pmod{p^{3r+3}}.
\]
Consequently all five $B$-record towers (T), including their nonlinear companions,
hold for every prime $p\ge5$ and every $r\ge2$.
\end{theorem}
The proof is given in Sections~\ref{sec:nu-new}--\ref{sec:completion}.
It eliminates the last arithmetic hypothesis in the five-record tower reduction.

\subsection{Method}
The starting point is that each summand factors exactly into an \emph{even} product
$E_k(n)=\prod_{i<k}(1-n^2/i^2)$, an \emph{odd} product $O_k(n)=\prod_{i<k}(1-n/i)$, and powers
of $x=n/k$ (Lemma~\ref{lem:reps}). At $n=p^r$ the factors with $p\mid i$ reproduce the
products at level $p^{r-1}$. The remaining factors, over reduced residues, are controlled by
two facts: power sums of reduced residues modulo $p^r$ vanish (Lemma~\ref{lem:units}), and
the reflection $i\mapsto p^r-i$ (Lemma~\ref{lem:units}(iii)). The second fact also yields the
comparison $\Delta_r\equiv 4p^r\sigma_1(p^r)$ with the reduced harmonic sum
(Lemma~\ref{lem:cb}).

Two telescoping identities (Lemma~\ref{lem:tel}) convert the $x^2$-weighted even sums, which
are not small, into $x^4$-weighted sums plus the boundary term $\binom{2n}n^2/4$. This boundary
term is what produces the constants $7/2$ and $-1/2$.

For the odd product the decisive input is the exact identity
$\sum_{k<n}(n/k)^2O_k(n)=nH_{n-1}$ for odd $n$ (Lemma~\ref{lem:alt}). It turns the
prime-level odd sum into a harmonic number, and that harmonic number is $\Delta_1/4$
modulo $p^5$.
For higher levels, polynomial moments of the zeroth-moment summands control
all complete residue blocks at once. Their recurrence loses at most
$v_p((j+1)!)$ at moment degree $j$; the block expansion supplies a compensating
factor $p^j$.

\subsection*{Conventions}
$p$ is an odd prime, $\vp$ the $p$-adic valuation on $\Q_p$, and all congruences are in $\Zp$.
For integers $X\ge1$ and $m\ge1$ let
\[
\sigma_m(X)=\sum_{\substack{1\le i<X\\ p\nmid i}}\frac1{i^m},\qquad H^{(m)}=H^{(m)}_{p-1}=\sigma_m(p),\qquad
H_n^{(m)}=\sum_{i=1}^n i^{-m},\quad H_n=H_n^{(1)}.
\]
For $1\le a,b$ we write $H(a,b)=\sum_{1\le i<k\le p-1}i^{-a}k^{-b}$.
When a level $r\ge1$ is fixed, $N=p^r$ and $M=p^{r-1}$.

\section{Exact identities}

For integers $n\ge1$ and $1\le k\le n$ set $x=x_{n,k}=n/k$ and
\[
O_k(n)=\prod_{i=1}^{k-1}\Bigl(1-\frac ni\Bigr),\quad
P_k(n)=\prod_{i=1}^{k-1}\Bigl(1+\frac ni\Bigr),\quad
E_k(n)=O_k(n)P_k(n)=\prod_{i=1}^{k-1}\Bigl(1-\frac{n^2}{i^2}\Bigr).
\]

\begin{lemma}\label{lem:fact}
For $1\le k\le n$:
\begin{align*}
\bin nk&=(-1)^{k-1}xO_k,& \bin{n+k}k&=(1+x)P_k,\\
\bin{n-1}k&=(-1)^k(1-x)O_k,&\bin{n+k-1}k&=xP_k .
\end{align*}
Moreover $O_n(n)=(-1)^{n-1}$, $P_n(n)=\tfrac12\binom{2n}n$ and $E_n(n)=(-1)^{n-1}\tfrac12\binom{2n}n$.
\end{lemma}
\begin{proof}
$\binom nk=\frac nk\prod_{i=1}^{k-1}\frac{n-i}{i}$, $\binom{n+k}k=\prod_{i=1}^k\frac{n+i}i$,
$\binom{n-1}k=\prod_{i=1}^k\frac{n-i}i$, and
$\binom{n+k-1}k=\prod_{i=1}^k\frac{n+i-1}i=\frac nk\prod_{i=1}^{k-1}\frac{n+i}i$.
The special values follow from the first two formulas at $k=n$.
\end{proof}

\begin{lemma}[representations]\label{lem:reps}
For $n\ge1$, with $x=n/k$ and $O_k=O_k(n)$, $E_k=E_k(n)$:
\begin{align*}
A(n)&=1+\sum_{k=1}^{n}x^2(1+x)^2E_k^2, & A(n-1)&=1+\sum_{k=1}^{n-1}x^2(1-x)^2E_k^2,\\
B(n)&=1+\sum_{k=1}^{n}x^2(1+x)O_kE_k, & B(n-1)&=1+\sum_{k=1}^{n-1}x(1-x)^2O_kE_k,\\
a(n)&=1-\sum_{k=1}^{n-1}x^3(1-x)E_k^2, & \nu(n)&=\sum_{k=1}^{n}O_kE_k .
\end{align*}
\end{lemma}
\begin{proof}
Insert Lemma~\ref{lem:fact} into each summand; for instance
$\binom{n-1}k^2\binom{n-1+k}{k}=(1-x)^2O_k^2\cdot xP_k=x(1-x)^2O_kE_k$, and
$\binom nk\binom{n-1}k\binom{n+k-1}k^2=-x^3(1-x)O_k^2P_k^2$. In $A(n-1)$ and $B(n-1)$ the term
$k=n$ vanishes. For $\nu$ put $k=j+1$ and use $\binom{n-1}{k-1}=(-1)^{k-1}O_k$ and
$\binom{n+k-1}{k-1}=P_k$.
\end{proof}

\begin{lemma}[telescoping]\label{lem:tel}
For $n\ge1$:
\[
\sum_{k=1}^{n-1}(x^4-2x^2)E_k^2=\tfrac14\bin{2n}n^2-1,\qquad
\sum_{k=1}^{n-1}(x^3-x^2-x)O_kE_k=\tfrac12\bin{2n}n-1 .
\]
\end{lemma}
\begin{proof}
$E_{k+1}^2-E_k^2=E_k^2\bigl((1-x^2)^2-1\bigr)$ and
$O_{k+1}E_{k+1}-O_kE_k=O_kE_k\bigl((1-x)(1-x^2)-1\bigr)$; sum over $k$ and use
Lemma~\ref{lem:fact} for $E_n^2$ and $O_nE_n=\tfrac12\binom{2n}n$.
\end{proof}

\begin{lemma}[alternating sums]\label{lem:alt}
For odd $n\ge1$:
\[
\sum_{k=1}^{n-1}xO_k=0,\qquad \sum_{k=1}^{n-1}x^2O_k=nH_{n-1},\qquad
\sum_{k=1}^{n-1}x^3O_k=\frac{n^2}2\bigl(H_n^2+H_n^{(2)}\bigr)-1 .
\]
\end{lemma}
\begin{proof}
By Lemma~\ref{lem:fact}, $x^{j}O_k=(-1)^{k-1}(n/k)^{j-1}\binom nk$. Put
$S_m(n)=\sum_{k=1}^n(-1)^{k-1}\binom nk k^{-m}$. Since
$\binom nk-\binom{n-1}k=\binom{n-1}{k-1}=\frac kn\binom nk$, we get
$S_m(n)-S_m(n-1)=S_{m-1}(n)/n$, with $S_0(n)=1$ for $n\ge1$. Hence $S_1(n)=H_n$ and
$S_2(n)=\sum_{j\le n}H_j/j=\tfrac12(H_n^2+H_n^{(2)})$. The term $k=n$ of $S_m(n)$ equals
$n^{-m}$ for odd $n$; subtracting it and multiplying by $n^{m}$ gives the three identities.
\end{proof}

Write
\begin{equation}\label{eq:IJ}
J_j(n)=\sum_{k=1}^{n-1}x^jE_k(n)^2,\qquad I_j(n)=\sum_{k=1}^{n-1}x^jO_k(n)E_k(n).
\end{equation}
Combining Lemmas~\ref{lem:reps} and~\ref{lem:tel}, and using that the term $k=n$ of
$A(n)$ is $4E_n^2=\binom{2n}n^2$ and that of $B(n)$ is $2O_nE_n=\binom{2n}n$, we obtain
with $C=\binom{2n}n$:
\begin{align}
A(n)&=\tfrac78C^2+\tfrac32+\tfrac32J_4+2J_3, &
A(n-1)&=-\tfrac18C^2+\tfrac32+\tfrac32J_4-2J_3, \label{eq:Arep}\\
a(n)&=1-J_3+J_4, &
B(n)&=1+C+I_2+I_3, \label{eq:arep}\\
B(n-1)&=2-\tfrac12C-3I_2+2I_3. \label{eq:Brep}
\end{align}
(For \eqref{eq:Arep} eliminate $J_2=\tfrac12(J_4-\tfrac14C^2+1)$; for \eqref{eq:Brep}
eliminate $I_1=I_3-I_2-\tfrac12C+1$.) In particular \eqref{eq:Arep}--\eqref{eq:arep} give
$12a(n)=A(n)+7A(n-1)$, confirming Bala's identity.

\section{\texorpdfstring{$p$}{p}-adic lemmas}

\begin{lemma}[reduced residues]\label{lem:units}
Let $p\ge5$, $r\ge1$, $N=p^r$, and $m\ge1$.
\begin{enumerate}
\item[(i)] If $(p-1)\nmid m$, then $\sigma_m(N)\equiv 0\pmod{p^r}$.
\item[(ii)] If $r\ge2$, then $\sigma_m(N)\equiv0\pmod p$.
\item[(iii)] (Reflection) If $m$ is odd, then
$2\sigma_m(N)=-\sum_{j\ge1}\binom{m+j-1}{j}N^j\sigma_{m+j}(N)$.
\item[(iv)] $\vp(\sigma_1(N))\ge2r$. If $p\ge7$, then $\vp(\sigma_3(N))\ge 2r$ and $\vp(\sigma_4(N))\ge r$.
For all $p\ge5$, one also has $\vp(\sigma_m(N))\ge r-1$ for every $m$,
and $\vp(\sigma_3(N))\ge2r-1$. If $r\ge2$, then $\vp(\sigma_3(N))\ge3$.
\end{enumerate}
\end{lemma}
\begin{proof}
(i) Let $g$ generate $(\Z/p^r\Z)^\times$. Modulo $p^r$ we have
$\sigma_m(N)\equiv\sum_{j=0}^{\varphi(N)-1}g^{-mj}=(g^{-m\varphi(N)}-1)/(g^{-m}-1)\equiv 0$,
because $g^{-m}-1$ is a unit when $(p-1)\nmid m$.
(ii) $i^{-m}\bmod p$ depends only on $i\bmod p$, and each nonzero class occurs $p^{r-1}$
times.
(iii) The map $i\mapsto N-i$ permutes the reduced residues in $[1,N)$, and
$(N-i)^{-m}=-i^{-m}(1-N/i)^{-m}=-\sum_{j\ge0}\binom{m+j-1}{j}N^ji^{-m-j}$ for odd $m$.
The series converges $p$-adically.
(iv) In (iii) with $m=1$, $\vp(N\sigma_2)\ge 2r$ by (i), and the other terms have valuation
$\ge2r$. For $m=3$ the same argument works when $\vp(\sigma_4)\ge r$, which (i) gives for
$p\ge7$. If $r\ge2$ then $\vp(N\sigma_4)\ge r+1\ge3$ by (ii) and $\vp(N^j\sigma_{3+j})\ge 4$
for $j\ge2$. To prove the general bound $v_p(\sigma_m(p^r))\ge r-1$,
split the reduced residues into lifts $a+tp^{r-1}$, $0\le t<p$.
For $r\ge2$, binomial expansion gives
$\sigma_m(p^r)\equiv p\sigma_m(p^{r-1})\pmod{p^{r-1}}$.
Induction proves the bound. Applying reflection to $\sigma_3$ then gives
$v_p(\sigma_3(N))\ge2r-1$, including $p=5$.
\end{proof}

\begin{lemma}[level $p$]\label{lem:wolst}
Let $p\ge7$. Then $\vp(H^{(1)})\ge2$, $\vp(H^{(2)})\ge1$, $\vp(H^{(3)})\ge2$, $\vp(H^{(4)})\ge1$,
\[
2H^{(1)}+pH^{(2)}\equiv0\pmod{p^4},\qquad H(2,2)\equiv H(1,3)\equiv0\pmod p .
\]
\end{lemma}
\begin{proof}
The four valuations are Lemma~\ref{lem:units} with $r=1$. By (iii) with $m=1$,
$2H^{(1)}+pH^{(2)}=-p^2H^{(3)}-p^3H^{(4)}-\sum_{j\ge4}p^jH^{(j+1)}$, and every term has
valuation $\ge4$. Next, $H(2,2)=\frac12\bigl((H^{(2)})^2-H^{(4)}\bigr)$. Finally
$H(1,3)+H(3,1)=H^{(1)}H^{(3)}-H^{(4)}\equiv0$, while the substitution
$(i,k)\mapsto(p-i,p-k)$ gives $H(1,3)\equiv H(3,1)\pmod p$.
\end{proof}

\begin{lemma}[central binomial]\label{lem:cb}
Let $p\ge5$ and $r\ge1$. Then
\begin{enumerate}
\item[(i)] $\binom{2N}N=\binom{2M}M\prod_{p\nmid i<N}(1+N/i)$;
\item[(ii)] $\vp(\Delta_r)\ge3r$, and $\binom{2p^s}{p^s}\equiv2\pmod{p^3}$ for all $s\ge0$;
\item[(iii)] if $p\ge7$, or if $p=5$ and $r\ge2$, then
$\Delta_r\equiv 4N\sigma_1(N)\pmod{p^{L_r}}$.
\end{enumerate}
Moreover $N\sigma_1(N)=NH_{N-1}-MH_{M-1}$ exactly.
\end{lemma}
\begin{proof}
(i) $\binom{2n}n=\prod_{i=1}^n(1+n/i)$. The factors with $i=pi'$, $1\le i'\le M$, equal
$1+M/i'$ and form $\binom{2M}M$.

(ii) Let $\Gamma=\prod_{p\nmid i<N}(1+N/i)$. Then
$\log\Gamma=\sum_{j\ge1}(-1)^{j-1}N^j\sigma_j(N)/j$. Here $\vp(N\sigma_1)\ge3r$ and
$\vp(N^2\sigma_2)\ge3r$ by Lemma~\ref{lem:units}, and $\vp(N^j\sigma_j/j)\ge jr-\vp(j)\ge 3r$
for $j\ge3$. Hence $\Gamma\equiv1\pmod{p^{3r}}$ and $\vp(\Delta_r)\ge3r$; induction on $s$ gives
the second claim.

(iii) By reflection,
$-\tfrac12N^2\sigma_2=N\sigma_1+\tfrac12\sum_{j\ge3}N^j\sigma_j$. Therefore
\[
\log\Gamma=2N\sigma_1+\sum_{j\ge3}\Bigl(\tfrac12+\tfrac{(-1)^{j-1}}j\Bigr)N^j\sigma_j .
\]
For $p\ge7$, $\vp(N^3\sigma_3)\ge5r$ and $\vp(N^4\sigma_4)\ge5r$ by Lemma~\ref{lem:units}(iv),
and the terms with $j\ge5$ have valuation $\ge 5r$. For $p=5,r\ge2$,
the corresponding lower bound is $5r-1$, using $v_p(\sigma_3)\ge2r-1$,
$v_p(\sigma_4)\ge r-1$, and $jr-v_p(j)\ge5r-1$ for $j\ge5$.
Both bounds are at least $L_r$, and $\vp(\log\Gamma)\ge3r$, so we
get $\Gamma-1\equiv2N\sigma_1\pmod{p^{L_r}}$. Then $\Delta_r=\binom{2M}M(\Gamma-1)$ together
with $\binom{2M}M\equiv2\pmod{p^3}$ (exactly $2$ if $r=1$) gives the claim.

The last identity holds because $H_{N-1}=\sigma_1(N)+H_{M-1}/p$.
\end{proof}

\begin{lemma}[carry sums]\label{lem:carry}
Let $p\ge5$, $K\ge1$ and $s=\vp(K)$. Then
\[
\sigma_m(pK)=\sum_{t\ge0}\bin{-m}{t}p^tF_t(K)H^{(m+t)},\qquad F_t(K)=\sum_{j=0}^{K-1}j^t,
\]
and the following hold:
\begin{enumerate}
\item[(i)] $\vp(\sigma_2(pK))\ge s+1$, and $\sigma_2(pK)\equiv KH^{(2)}\pmod{p^{s+2}}$;
\item[(ii)] if $p\ge7$, then $\sigma_2(pK)\equiv KH^{(2)}\pmod{p^{s+3}}$;
\item[(iii)] $\vp(\sigma_m(pK))\ge\max(0,s-1)$ for all $m$;
\item[(iv)] $\vp(\sigma_1(pK))\ge2s+2$.
\end{enumerate}
\end{lemma}
\begin{proof}
Write $i=jp+u$ with $0\le j<K$, $1\le u\le p-1$, and expand $(jp+u)^{-m}$ binomially.
By Faulhaber's formula and von Staudt--Clausen, $\vp(F_t(K))\ge s-1-\vp(t+1)$. The explicit
forms $F_1=\frac{K(K-1)}2$ and $F_2=\frac{K(K-1)(2K-1)}6$ give $\vp(F_1),\vp(F_2)\ge s$, and
$F_3=\frac{K^2(K-1)^2}{4}$ gives $\vp(F_3)\ge 2s$.

For $m=2$ the $t$-th term is $(-1)^t(t+1)p^tF_tH^{(t+2)}$. The term $t=0$ is $KH^{(2)}$, of
valuation $\ge s+1$. For $t=1$ the valuation is $\ge 1+s+\vp(H^{(3)})\ge s+2$, and for $t=2$
it is $\ge 2+s$. Both are $\ge s+3$ when $p\ge7$, because then $\vp(H^{(3)})\ge2$ and
$\vp(H^{(4)})\ge1$. For $t=3$ it is $\ge 2s+3$. For
$t\ge4$ it is $\ge t+s-1-\vp(t+1)$, which is $\ge s+2$, and $\ge s+3$ when $p\ge7$.

(iii) follows from the same bounds, together with integrality.

(iv) Reflection on the reduced residues in $[1,pK)$ gives
$2\sigma_1(pK)=-\sum_{j\ge1}(pK)^j\sigma_{1+j}(pK)$. The term $j=1$ has valuation
$\ge1+s+(s+1)$ by (i), and the terms $j\ge2$ have valuation $\ge 2(1+s)$ by (iii).
\end{proof}

\begin{lemma}[carry factors]\label{lem:cf}
Let $p\ge5$, $r\ge2$, $1\le K\le M-1$ and $s=\vp(K)$, so that $s\le r-2$. Put
\[
\mathcal A_K=\prod_{p\nmid i<pK}\Bigl(1-\frac Ni\Bigr),\qquad
\mathcal B_K=\prod_{p\nmid i<pK}\Bigl(1-\frac{N^2}{i^2}\Bigr).
\]
Then $\mathcal B_K-1+N^2\sigma_2(pK)$ and $\mathcal B_K^2-1+2N^2\sigma_2(pK)$ have valuation
$\ge 4r+s-1$, $\vp(\mathcal B_K-1)\ge 2r+s+1$, and
$\vp(\mathcal A_K-1)\ge\min(r+2s+2,\ 2r+s+1)$. Moreover
$O_{pK}(N)=O_K(M)\mathcal A_K$ and $E_{pK}(N)=E_K(M)\mathcal B_K$.
\end{lemma}
\begin{proof}
For the factorisations, the factors with $i=pi'$ are $1-M/i'$ and $1-M^2/i'^2$. For the
estimates, $\log\mathcal B_K=-\sum_{j\ge1}N^{2j}\sigma_{2j}(pK)/j$. The terms with $j\ge2$
have valuation $\ge 4r+s-1$ by Lemma~\ref{lem:carry}(iii), and the $j=1$ term has valuation
$\ge2r+s+1$ by Lemma~\ref{lem:carry}(i). Exponentiating gives the statements about
$\mathcal B_K$. Similarly $\log\mathcal A_K=-\sum_jN^j\sigma_j(pK)/j$, where the $j=1$ term has
valuation $\ge r+2s+2$ by Lemma~\ref{lem:carry}(iv), the $j=2$ term $\ge 2r+s+1$, and the terms
with $j\ge3$ have valuation $\ge 3r-\vp(j)$, which exceeds both bounds.
\end{proof}

\begin{lemma}[top level]\label{lem:top}
Let $p\ge5$, $r\ge2$, $p\nmid k<N$, and put $q=\lfloor k/p^{r-1}\rfloor\in\{0,\dots,p-1\}$ and
$h_q=\sum_{u=1}^qu^{-2}$. Then $O_k(N)\equiv1\pmod p$,
$E_k(N)\equiv1-p^2h_q\pmod{p^4}$ and $E_k(N)^2\equiv1-2p^2h_q\pmod{p^4}$. Moreover, for
each $q$,
\[
\beta_q:=\sum_{\substack{p\nmid k\\ qp^{r-1}<k<(q+1)p^{r-1}}}k^{-3}\equiv0\pmod p .
\]
\end{lemma}
\begin{proof}
For $i<k$ we have $\vp(N^2/i^2)=2(r-\vp(i))\ge2$, with equality exactly for $i=p^{r-1}u$,
$1\le u\le q$. Products of two deviations have valuation $\ge4$, and $\vp(N/i)\ge1$.

For $\beta_q$: the block contains every nonzero residue class modulo $p$ exactly $p^{r-2}$
times, so $\beta_q\equiv p^{r-2}H^{(3)}\equiv 0\pmod p$, since $(p-1)\nmid3$.
\end{proof}

\section{Proof of Theorem~\ref{thm:A212334}}

\subsection*{\texorpdfstring{Level $r=1$}{Level r=1}}
By Lemma~\ref{lem:reps}, $a(p)-1=-\sum_{k=1}^{p-1}(p/k)^3(1-p/k)E_k(p)^2$, and
$E_k(p)^2\equiv1\pmod{p^2}$. Hence $a(p)-1\equiv-p^3H^{(3)}+p^4H^{(4)}\equiv0\pmod{p^5}$ for
$p\ge7$ by Lemma~\ref{lem:wolst}. For $p=5$ one computes $a(5)-1=87500=28\cdot5^5$.

\subsection*{\texorpdfstring{Levels $r\ge2$}{Levels r at least 2}}
Put $\tau(n,k)=x^3(1-x)E_k(n)^2$, so $a(n)=1-\sum_{k<n}\tau(n,k)$. Splitting $1\le k<N$
according to whether $p\mid k$ gives
\[
a(N)-a(M)=-(U+V),\qquad U=\sum_{p\nmid k<N}\tau(N,k),\qquad
V=\sum_{K=1}^{M-1}\bigl(\tau(N,pK)-\tau(M,K)\bigr).
\]

\emph{The unit part.} For $p\nmid k$ we have $\vp(N/k)=r\ge2$, and by Lemma~\ref{lem:top},
$(1-N/k)E_k(N)^2=1-N/k-2p^2h_q+O(p^4)$. Hence
\[
U\equiv N^3\Bigl(\sigma_3(N)-N\sigma_4(N)-2p^2\sum_{q=0}^{p-1}h_q\beta_q\Bigr)\pmod{p^{3r+4}}.
\]
Here $\sigma_3(N)\equiv0\pmod{p^3}$ by Lemma~\ref{lem:units}(iv), $\vp(N\sigma_4(N))\ge r+1\ge3$
by Lemma~\ref{lem:units}(ii), and $\beta_q\equiv0\pmod p$. Thus $U\equiv0\pmod{p^{3r+3}}$.

\emph{The carry part.} Since $N/(pK)=M/K$ and $E_{pK}(N)=E_K(M)\mathcal B_K$, we have
$\tau(N,pK)=\tau(M,K)\mathcal B_K^2$, and therefore
$V=\sum_K\tau(M,K)(\mathcal B_K^2-1)$. Let $s=\vp(K)\le r-2$. Then
$\vp(\tau(M,K))=3(r-1-s)$, because $1-M/K$ and $E_K(M)$ are units. By Lemma~\ref{lem:cf},
\[
\tau(M,K)(\mathcal B_K^2-1)=-2N^2\tau(M,K)\sigma_2(pK)+O\bigl(p^{3(r-1-s)+4r+s-1}\bigr).
\]
The error has valuation $\ge 7r-4-2s\ge 5r\ge3r+3$. The main term has valuation
$\ge5r-2-2s$, which is $\ge 3r+4$ unless $s=r-2$. If $s=r-2$ then $K=p^{r-2}m$ with
$1\le m\le p-1$. In that case $\tau(M,K)\equiv p^3m^{-3}\pmod{p^4}$ and
$\sigma_2(pK)\equiv p^{r-2}mH^{(2)}\pmod{p^{r}}$ by Lemma~\ref{lem:carry}(i), so
\[
\sum_{m=1}^{p-1}-2N^2\tau(M,p^{r-2}m)\sigma_2(p^{r-1}m)\equiv
-2p^{3r+1}H^{(2)}\sum_{m=1}^{p-1}m^{-2}=-2p^{3r+1}\bigl(H^{(2)}\bigr)^2\equiv0\pmod{p^{3r+3}}.
\]
Hence $V\equiv0$, and $a(N)\equiv a(M)\pmod{p^{3r+3}}$. \qed

\section{Proof of Theorem~\ref{thm:W}}

\begin{lemma}\label{lem:J4}
Under the hypotheses of Theorem~\ref{thm:W}, $J_4(N)\equiv J_4(M)\pmod{p^{L_r}}$.
\end{lemma}
\begin{proof}
For $r=1$: $J_4(1)=0$ and $J_4(p)\equiv p^4H^{(4)}\equiv0\pmod{p^5}$ for $p\ge7$.

For $r\ge2$ split as before. The unit part $N^4\sum_{p\nmid k}k^{-4}E_k(N)^2$ has valuation
$\ge4r\ge3r+3$ when $r\ge3$. When $r=2$ it is $\equiv N^4\sigma_4(p^2)\equiv0\pmod{p^{9}}$ by
Lemma~\ref{lem:units}(ii). The carry part is $\sum_K(M/K)^4E_K(M)^2(\mathcal B_K^2-1)$, and its
terms have valuation $\ge 4(r-1-s)+2r+s+1=6r-3-3s\ge3r+3$.
\end{proof}

\begin{proof}[Proof of Theorem~\ref{thm:W}]
By \eqref{eq:arep}, $J_3(N)-J_3(M)=\bigl(J_4(N)-J_4(M)\bigr)-\bigl(a(N)-a(M)\bigr)\equiv0$, using
Lemma~\ref{lem:J4} and Theorem~\ref{thm:A212334}. By Lemma~\ref{lem:cb}(ii),
\[
\bin{2N}N^2-\bin{2M}M^2=\Delta_r\Bigl(\bin{2N}N+\bin{2M}M\Bigr)\equiv4\Delta_r\pmod{p^{3r+3}}.
\]
Taking the differences of \eqref{eq:Arep} at $n=N$ and $n=M$ gives
$\gamma_r\equiv\frac78\cdot4\Delta_r$ and $\delta_r\equiv-\frac18\cdot4\Delta_r$ modulo
$p^{L_r}$.
\end{proof}

\begin{remark}
For $p=5$ and $r=1$ the statement of Theorem~\ref{thm:W} fails. One has
$\vp\bigl(A(5)-5-\tfrac72\Delta_1\bigr)=4$, because $(p-1)\mid4$ makes $H^{(4)}$ a unit.
Bala's combinations are nevertheless still true at $p=5$; see Corollary~\ref{cor:prime}.
\end{remark}

\section{Prime level for \texorpdfstring{$B$}{B}: Theorem~\ref{thm:Z1}, Corollary~\ref{cor:prime}, Theorem~\ref{thm:nu}}

\begin{proof}[Proof of Theorem~\ref{thm:Z1}]
Let $p\ge7$. For $1\le k\le p-1$ we have $O_k(p)\equiv1-pH_{k-1}\pmod{p^2}$ and
$E_k(p)\equiv1-p^2H^{(2)}_{k-1}\pmod{p^4}$.

First, $I_3(p)\equiv p^3H^{(3)}-p^4H(1,3)\equiv0\pmod{p^5}$ by Lemma~\ref{lem:wolst}.

Second, by Lemma~\ref{lem:alt},
\[
I_2(p)=\sum_{k<p}x^2O_k+\sum_{k<p}x^2O_k(E_k-1)
\equiv pH_{p-1}-p^4H(2,2)\equiv pH_{p-1}\pmod{p^5}.
\]
By Lemma~\ref{lem:cb}(iii) with $r=1$, $pH_{p-1}=p\sigma_1(p)\equiv\frac14\Delta_1$.

Now \eqref{eq:arep}--\eqref{eq:Brep} with $\binom{2p}p=2+\Delta_1$ give
\[
B(p)\equiv 3+\Delta_1+\tfrac14\Delta_1,\qquad
B(p-1)\equiv 2-\tfrac12(2+\Delta_1)-\tfrac34\Delta_1=1-\tfrac54\Delta_1\pmod{p^5}.\qedhere
\]
\end{proof}

\begin{proof}[Proof of Corollary~\ref{cor:prime}]
Let $p\ge7$ and $\varepsilon=\Delta_1$, so $\vp(\varepsilon)\ge3$ and $\varepsilon^2\equiv0\pmod{p^6}$.
By Theorems~\ref{thm:W} and~\ref{thm:Z1}, modulo $p^5$,
\[
A(p)\equiv5+\tfrac72\varepsilon,\quad A(p-1)\equiv1-\tfrac12\varepsilon,\quad
B(p)\equiv3+\tfrac54\varepsilon,\quad B(p-1)\equiv1-\tfrac54\varepsilon .
\]
For the linear statements, the coefficient of $\varepsilon$ vanishes in each case:
\[
\tfrac12\bigl(\tfrac54-\tfrac54\bigr)=0,\quad
5\cdot\tfrac72-14\cdot\tfrac54=0,\quad
5\cdot\bigl(-\tfrac12\bigr)-2\cdot\bigl(-\tfrac54\bigr)=0,\quad
5\cdot\bigl(-\tfrac12\bigr)+2\cdot\tfrac54=0 .
\]
The constants $2$, $-17$, $3$ and $11$ are the values at $p^0=1$.

For the nonlinear statements, expand to first order in $\varepsilon$:
\begin{align*}
A(p)^{25}&\equiv5^{25}\bigl(1+\tfrac{35}2\varepsilon\bigr), &
B(p)^{42}&\equiv3^{42}\bigl(1+\tfrac{35}2\varepsilon\bigr),\\
A(p-1)^5B(p)^6&\equiv\bigl(1-\tfrac52\varepsilon\bigr)3^6\bigl(1+\tfrac52\varepsilon\bigr)\equiv3^6, &
B(p)^3B(p-1)&\equiv27\bigl(1+\tfrac54\varepsilon\bigr)\bigl(1-\tfrac54\varepsilon\bigr)\equiv27 .
\end{align*}

For $p=5$ the statements are finite computations. With $A(5)=819005$, $A(4)=33001$,
$B(5)=11253$ and $B(4)=1251$, the six differences have $5$-adic valuations
$5,6,5,6,\ge25,6$ in the order displayed in Corollary~\ref{cor:prime}, and
$B(5)^3B(4)-27$ has valuation $5$. For $p=3$, $B(3)^3B(2)-27=147^3\cdot19-27=3^6\cdot 82790$.
\end{proof}

\begin{proof}[Proof of Theorem~\ref{thm:nu}]
By Lemma~\ref{lem:reps}, $\nu(p)=\sum_{k=1}^pO_k(p)E_k(p)$. By Lemma~\ref{lem:fact},
$\sum_{k=1}^pO_k(p)=\sum_{j=0}^{p-1}(-1)^j\binom{p-1}j=0$, so
$\nu(p)=\sum_{k=1}^{p}O_k(E_k-1)$. For $k\le p$,
$O_k(E_k-1)\equiv -p^2H^{(2)}_{k-1}+p^3H_{k-1}H^{(2)}_{k-1}\pmod{p^4}$. Now
\[
\sum_{m=0}^{p-1}H^{(2)}_m=pH^{(2)}-H^{(1)},
\]
which has valuation $\ge2$ for $p\ge5$. Moreover $H_{p-1-m}\equiv H_m$ and
$H^{(2)}_{p-1-m}\equiv-H^{(2)}_m\pmod p$, so the map $m\mapsto p-1-m$ shows
$\sum_{m=0}^{p-1}H_mH^{(2)}_m\equiv0\pmod p$. Hence $\nu(p)\equiv0\pmod{p^4}$.
\end{proof}

\section{Towers for \texorpdfstring{$B$}{B}: proof of Theorem~\ref{thm:towers}}

(a) is Lemma~\ref{lem:cb}(iii) together with the identity at the end of that lemma.

(b) Write $I_3=F_3+R_3$ with $F_3(n)=\sum_{k<n}x^3O_k$ and
$R_3(n)=\sum_{k<n}x^3O_k(E_k-1)$.

\emph{The term $F_3$.} By Lemma~\ref{lem:alt},
$F_3(n)=\frac{n^2}{2}(H_n^2+H_n^{(2)})-1$. Put $\sigma_j=\sigma_j(N)$. Using
$H_N=\sigma_1+H_M/p$ and $H_N^{(2)}=\sigma_2+H_M^{(2)}/p^2$,
\[
F_3(N)-F_3(M)=\tfrac12\bigl(N^2\sigma_1^2+2N\sigma_1\cdot MH_M+N^2\sigma_2\bigr).
\]
Here $MH_M=1+MH_{M-1}$ and $MH_{M-1}=\sum_{t=1}^{r-1}p^t\sigma_1(p^t)\equiv0\pmod{p^3}$, the
latter by telescoping and Lemma~\ref{lem:units}(iv). Hence
$F_3(N)-F_3(M)\equiv N\sigma_1+\frac12N^2\sigma_2\pmod{p^{3r+3}}$. By reflection this equals
$-\frac12\sum_{j\ge2}N^{j+1}\sigma_{j+1}$, which has valuation $\ge5r\ge L_r$ for $p\ge7$,
and $\ge5r-1\ge L_r$ for $p=5,r\ge2$.

\emph{The term $R_3$.} For $r=1$, $\vp(R_3(p))\ge3+2$. For $r\ge2$ split into unit and carry
parts.

The unit part is $\sum_{p\nmid k}(N/k)^3O_k(N)(E_k(N)-1)$. By Lemma~\ref{lem:top} it is
$\equiv-p^2N^3\sum_qh_q\beta_q\equiv0\pmod{p^{3r+3}}$.

For the carry part, Lemma~\ref{lem:cf} gives
\[
O_{pK}(N)(E_{pK}(N)-1)-O_K(M)(E_K(M)-1)=O_K(M)\bigl[(E_K(M)-1)(\mathcal A_K\mathcal B_K-1)+\mathcal A_K(\mathcal B_K-1)\bigr].
\]
Multiply by $(M/K)^3$, which has valuation $3(r-1-s)$. The first summand then has valuation
$\ge\min(4r+1-s,\ 5r-2s)\ge3r+3$. The second has valuation $\ge5r-2-2s$, which is
$\ge3r+4$ unless $s=r-2$. For $K=p^{r-2}m$ it is $\equiv-p^{3r+1}H^{(2)}m^{-2}$, exactly as in
\S4, and the sum over $m$ vanishes modulo $p^{3r+3}$.

(c) By \eqref{eq:arep}--\eqref{eq:Brep},
$\alpha_r=\Delta_r+\Delta I_2+\Delta I_3$ and
$\beta_r=-\frac12\Delta_r-3\Delta I_2+2\Delta I_3$, where $\Delta I_j=I_j(N)-I_j(M)$.
Insert (b) and the definition of $\mathcal D_r$, then use Theorem~\ref{thm:W}:
\[
\alpha_r+\beta_r\equiv-2\mathcal D_r,\qquad 5\gamma_r-14\alpha_r\equiv-14\mathcal D_r,\qquad
5\delta_r-2\beta_r\equiv6\mathcal D_r\pmod{p^{L_r}}.\qquad\qed
\]

% Independently checked addition to apery_rigidity.tex.
\section{The zeroth moment at all prime powers}\label{sec:nu-new}

\begin{proposition}\label{prop:nu-exact}
For every integer $n\geq 1$,
\[
 5(n-1)\nu(n)=n\bigl(B(n)-3B(n-1)\bigr).
\]
\end{proposition}
\begin{proof}
Put $t_k=\binom nk^2\binom{n+k}k$ for $0\leq k\leq n$, and set
$t_{n+1}=0$. The summands of $\nu(n)$ and $B(n-1)$, extended by a
zero term at $k=n$, are respectively
\[
 \frac{(n-k)^2}{n^2}t_k,
 \qquad \frac{(n-k)^2}{n(n+k)}t_k.
\]
Define
\[
 f(k)=\frac{5(n-1)(n-k)^2}{n^2}-n+
             \frac{3(n-k)^2}{n+k},
 \qquad G(k)=\frac{k^3(7n-5k)}{n^2(n+k)}.
\]
The quotient $t_{k+1}/t_k=(n-k)^2(n+k+1)/(k+1)^3$ and the polynomial
identity
\[
 (n-k)^2(n+k)(7n-5k-5)-k^3(7n-5k)
   =n^2(n+k)f(k)
\]
give $G(k+1)t_{k+1}-G(k)t_k=f(k)t_k$ for $0\leq k\leq n$.
Both boundary terms vanish, because $G(0)=0$ and $t_{n+1}=0$.
Summing proves the assertion, including $n=1$.
\end{proof}


\begin{theorem}\label{thm:nu-all-powers}
For every prime $p\ge5$ and every $s\ge1$,
\[
v_p(\nu(p^s))\ge s+3.
\]
\end{theorem}
\begin{proof}
Substitution of the exact $B$ representations into
Proposition~\ref{prop:nu-exact} cancels the factor $5$ before any reduction:
for $n>1$,
\[
\nu(n)=\frac{n}{n-1}
\left(\frac12\binom{2n}{n}-1+2I_2(n)-I_3(n)\right).
\]
For $n=p^s$, the first two terms in parentheses have difference in
$p^3\mathbb Z_p$ by Lemma~\ref{lem:cb}.
Also $I_3(n)\in p^3\mathbb Z_p$ termwise. The alternating identity gives
\[
I_2(n)=nH_{n-1}+R(n),\qquad
R(n)=\sum_{k<n}(n/k)^2O_k(n)(E_k(n)-1).
\]
Here $R(n)\in p^4\mathbb Z_p$ termwise, since
$v_p(n/k)\ge1$ and $E_k(n)-1\in p^2\mathbb Z_p$.
Finally,
$nH_{n-1}=\sum_{a=1}^s p^a\sigma_1(p^a)\in p^3\mathbb Z_p$.
Thus the entire parenthesis is in $p^3\mathbb Z_p$; multiplication by
$n/(n-1)$ proves the assertion. This proof includes $p=5$ without division
by a nonunit and uses no unproved enhanced congruence.
\end{proof}

% Standalone proof fragment; notation agrees with apery_rigidity.tex.
\section{The unit contribution: a moment argument}

We use the previously established bound
\begin{equation}\label{eq:nu-input-unit}
 v_p(\nu(p^s))\ge s+3\qquad(p\ge5,\ s\ge1).
\end{equation}
All the arguments below are finite identities or convergent $p$-adic series.

\begin{lemma}[Moment divisibility]\label{lem:moment-unit}
Let $p\ge5$, $s\ge1$, $n=p^s$, and put
\begin{gather*}
Q_l=O_{l+1}(n)=(-1)^l\binom{n-1}{l},\qquad T_l=E_{l+1}(n),\\
W_l=Q_lT_l=\binom{n-1}{l}^{\!2}\binom{n+l}{l}\qquad(0\le l<n).
\end{gather*}
Write $S_j=\sum_{l=0}^{n-1}l^jW_l$ and
$P_j=\sum_{l=0}^{n-1}l^jQ_l(T_l-1)$, with $l^0=1$ even at $l=0$.
Let $t=\min(2s,s+3)$. Then $S_0=\nu(n)$, and
\begin{align}
 (2-n)S_1&=(2n-1)\nu(n)-n^2B(n-1),\label{eq:moment-first}\\
 (n-3)S_2+(n^2+2n-3)S_1-(n-1)^2(n+1)\nu(n)&=0.\label{eq:moment-second}
\end{align}
For $1\le j<n-1$,
\[
 P_j=S_j,\qquad v_p(S_j)\ge t-v_p((j+1)!).
\]
Also $P_0=\nu(n)$, $v_p(P_j)\ge2$ for every $j\ge0$, and
$v_p(S_1+S_2)\ge s+2$.
\end{lemma}
\begin{proof}
The alternating binomial finite-difference identity gives
$\sum_l l^jQ_l=0$ for $0\le j<n-1$, proving $P_j=S_j$ in this range.
Every factor of $T_l$ is $1-n^2/i^2$ with $v_p(n/i)\ge1$, so
$T_l-1\in p^2\mathbb Z_p$; hence $P_j\in p^2\mathbb Z_p$ for all $j$.

For \eqref{eq:moment-first}, put
$b_l=\binom{n-1}{l}^2\binom{n-1+l}{l}$ and $b_n=0$.
Then $W_l=(n+l)b_l/n$, whence
$\sum_l b_l=B(n-1)$ and $\sum_l l b_l=n(\nu(n)-B(n-1))$.
The relation
\[
(l+1)^3b_{l+1}=(n-1-l)^2(n+l)b_l
\]
can be summed after subtracting $l^3b_l$. It gives
\[
(2-n)\sum_l l^2b_l+(1-n^2)\sum_l l b_l+n(n-1)^2\sum_l b_l=0.
\]
Combining these equations yields \eqref{eq:moment-first}; the factor $n$ is
cancelled in an exact rational identity, before taking valuations. Since
$2-n$ is a $p$-adic unit, $v_p(S_1)\ge t$.

For $j\ge2$, use $W_n=0$ and
\[
(l+1)^3W_{l+1}=(l+1-n)^2(l+1+n)W_l.
\]
Multiplication by $(l+1)^{j-2}$ and summation gives the exact identity
\begin{equation}\label{eq:moment-polynomial}
\sum_{l=0}^{n-1}
\left[l^{j+1}-(l+1)^{j-2}(l+1-n)^2(l+1+n)\right]W_l=0.
\end{equation}
The bracket is an integer polynomial in $l$ of degree $j$, with leading
coefficient $n-j-1$. The case $j=2$ gives \eqref{eq:moment-second}.
When $j+1<n=p^s$, we have
$v_p(n-j-1)=v_p(j+1)$. Induction in \eqref{eq:moment-polynomial}, beginning
with $S_0\in p^{s+3}\mathbb Z_p$ and $S_1\in p^t\mathbb Z_p$, gives
$v_p(S_j)\ge t-v_p((j+1)!)$.
Finally, \eqref{eq:moment-second} gives
\[
(n-3)(S_1+S_2)=-n(n+1)S_1+(n-1)^2(n+1)\nu(n).
\]
Here $n-3$ is a unit and $s+t\ge s+2$, so the last assertion follows.
\end{proof}

\begin{lemma}[Weighted complete blocks]\label{lem:weighted-complete}
Under the notation and assumptions of Lemma~\ref{lem:moment-unit},
\[
\sum_{l=0}^{n-1}Q_l(T_l-1)
 \sum_{u=1}^{p-1}\frac1{(pl+u)^2}\equiv0\pmod{p^{s+4}}.
\]
\end{lemma}
\begin{proof}
Expanding each denominator gives the convergent series
\[
\sum_{j\ge0}(-1)^j(j+1)p^jH^{(j+2)}P_j.
\]
The $j=0$ term has valuation at least $1+(s+3)=s+4$.
If $p\ge7$, the coefficients of $P_1$ and $P_2$ have valuation at least
$3$, by $v_p(H^{(3)})\ge2$ and $v_p(H^{(4)})\ge1$. Since
$v_p(S_1),v_p(S_2)\ge t\ge s+1$, these terms vanish at the desired precision.
For $p=5$, reflection gives
$H^{(3)}\equiv-\tfrac32pH^{(4)}\pmod{p^2}$.
Consequently the sum of these two terms is
\[
-2pH^{(3)}S_1+3p^2H^{(4)}S_2
\equiv3p^2H^{(4)}(S_1+S_2)\pmod{p^{s+4}},
\]
because the discarded coefficient has valuation at least $3$ and
$v_p(S_1)\ge s+1$. The surviving expression vanishes by
$v_p(S_1+S_2)\ge s+2$.

For $3\le j<n-1$, use
\[
j-v_p((j+1)!)\ge3\qquad(p\ge5,\ j\ge3).
\]
For $j=3,4$ this follows directly, and for $j\ge5$ it follows from
$v_p((j+1)!)\le(j+1)/(p-1)\le(j+1)/4$.
Thus the $j$th term has valuation at least $t+3\ge s+4$.
For $j\ge n-1$, the elementary bound $v_p(P_j)\ge2$ gives valuation
at least $j+2\ge n+1\ge s+4$. This proves the assertion, including the
infinite tail.
\end{proof}

\begin{theorem}[The entire unit contribution]\label{thm:unit-all}
Let $p\ge5$, $r\ge2$, $N=p^r$, and $M=p^{r-1}$. Then
\[
 U_2=\sum_{\substack{1\le k<N\\p\nmid k}}
 \left(\frac Nk\right)^2O_k(N)(E_k(N)-1)
 \equiv0\pmod{p^{3r+3}}.
\]
\end{theorem}
\begin{proof}
Write $k=pl+u$, $0\le l<M$, $1\le u<p$, and take $Q_l,T_l$ from
Lemma~\ref{lem:moment-unit} with $n=M$ and $s=r-1$.
Define the products over the nonmultiples of $p$ below $k$ by
\[
 A_k=\prod_{\substack{i<k\\p\nmid i}}(1-N/i),\qquad
 B_k^*=\prod_{\substack{i<k\\p\nmid i}}(1-N^2/i^2),
\]
and the corresponding partial sums by
$\sigma_a(k)=\sum_{i<k,\,p\nmid i}i^{-a}$.
Then $O_k(N)=Q_lA_k$ and $E_k(N)=T_lB_k^*$ exactly. Put
$D_l=Q_l(T_l-1)$. Since $v_p(D_l)\ge2$, product expansion gives, termwise,
\begin{equation}\label{eq:unit-expand}
O_k(N)(E_k(N)-1)
\equiv D_l-ND_l\sigma_1(k)-N^2Q_lT_l\sigma_2(k)
\pmod{p^{r+3}}.
\end{equation}
Indeed, the quadratic remainder of $A_k$ multiplied by $D_l$ has
valuation at least $2r+2\ge r+3$; the remaining cross terms have
valuation at least $3r\ge r+3$; and the quadratic remainder of $B_k^*$
has valuation at least $4r\ge r+3$.

The contribution of $D_l$ to $U_2/N^2$ is zero modulo $p^{r+3}$ by
Lemma~\ref{lem:weighted-complete}.
For the term with $\sigma_1(k)$, it suffices to prove
\[
\sum_{l,u}D_l\sigma_1(pl+u)(pl+u)^{-2}\equiv0\pmod{p^3}.
\]
Let $q=\lfloor l/p^{r-2}\rfloor$ and $h_q=\sum_{a=1}^qa^{-2}$.
Only the factors whose indices have valuation $r-2$ contribute to
$T_l-1$ modulo $p^3$, so $D_l/p^2\equiv-h_q\pmod p$.
Furthermore,
$\sigma_1(pl+u)\equiv H_{u-1}\pmod p$ and
$(pl+u)^{-2}\equiv u^{-2}\pmod p$. Hence the displayed sum divided by
$p^2$ is congruent to
\[
-H(1,2)\,p^{r-2}\sum_{q=0}^{p-1}h_q\pmod p.
\]
This is zero: indeed
$\sum_{q=0}^{p-1}h_q=pH^{(2)}-H^{(1)}\equiv0\pmod p$.
This argument works also when $r=2$.

The term with $N^2\sigma_2(k)$ in \eqref{eq:unit-expand} is already
zero modulo $p^{r+3}$ for $r\ge3$. If $r=2$, it suffices to consider its
remaining sum modulo $p$. We have $Q_lT_l\equiv1\pmod p$ and
\[
\sigma_2(pl+u)\equiv H_{u-1}^{(2)}\pmod p,
\]
since each complete block contributes $H^{(2)}\equiv0\pmod p$.
There are exactly $M=p$ choices of $l$, and consequently
\[
\sum_{l=0}^{p-1}\sum_{u=1}^{p-1}
 Q_lT_l\sigma_2(pl+u)(pl+u)^{-2}
\equiv p\sum_{u=1}^{p-1}u^{-2}H_{u-1}^{(2)}\equiv0\pmod p.
\]
All contributions to $U_2/N^2$ therefore vanish modulo $p^{r+3}$,
which proves the theorem.
\end{proof}

% Intermediate proof, for integration into the revised manuscript.
\section{The carry contribution}\label{sec:carry-new}
\begin{lemma}[The carry contribution]\label{lem:carry-new}
Let $p\ge5$, $r\ge2$, $N=p^r$, $M=p^{r-1}$, and use the notation of the manuscript.
If $\nu(M)\equiv0\pmod{p^3}$, then the carry part of
$R(N)-R(M)$ vanishes modulo $p^{3r+3}$.
\end{lemma}
\begin{proof}
Put $x=M/K$, $s=v_p(K)$, $O=O_K(M)$, $E=E_K(M)$,
$A=\mathcal A_K$, and $B=\mathcal B_K$.
The carry contribution is exactly
\[
 V=\sum_{K=1}^{M-1}x^2O\{(E-1)(AB-1)+A(B-1)\}.
\]
We use $v_p(x)=r-1-s$, $s\le r-2$, $v_p(E-1)\ge2$,
$v_p(A-1)\ge r+2s+2$, and $v_p(B-1)\ge2r+s+1$.
Write $H_j=H_{p-1}^{(j)}$.
The carry harmonic sums satisfy, for all $p\ge5$,
\begin{align}
 \sigma_2(pK)&\equiv K H_2+p^2(K^3-K)H_4\pmod{p^{s+3}},\
 \sigma_1(pK)&\equiv-\tfrac12pK^2H_2\pmod{p^{2s+3}}.
\end{align}
For $p\ge7$, the $H_4$ correction in the first line is zero at the displayed precision.
Indeed, expand the carry sum by blocks. The terms of orders zero, one, and two are
$KH_2-2pF_1(K)H_3+3p^2F_2(K)H_4$.
Reflection gives $H_3\equiv-3pH_4/2\pmod{p^2}$, and
$3(F_1+F_2)=K^3-K$.
The term of order three has valuation at least $s+3$, since
$F_3=K^2(K-1)^2/4$; for orders $t\ge4$, the estimate
$v_p(F_t)\ge s-1-v_p(t+1)$ and the factor $(t+1)p^t$ give the same conclusion.
For the second line use reflection with $T=pK$:
$2\sigma_1(T)=-T\sigma_2(T)-T^2\sigma_3(T)-T^3\sigma_4(T)-\cdots$.
Here $v_p(\sigma_4(T))\ge s$ and $v_p(\sigma_3(T))\ge2s+1$.
The former follows directly from the same block expansion, and the latter follows by reflection.
Together with $\sigma_2(T)\equiv KH_2\pmod{p^{s+2}}$, these give the claimed precision.

Consequently, in the weighted sum defining $V$, we may replace
\[
 A-1\quad\hbox{by}\quad a_0=\tfrac12(NpK^2-N^2K)H_2,
\]
and replace
\[
 B-1\quad\hbox{by}\quad b_0+c_0,
 \qquad b_0=-N^2KH_2,\quad c_0=-N^2p^2(K^3-K)H_4.
\]
For clarity, the error in the first replacement is multiplied by $x^2(E-1)$.
Its $N\sigma_1$ error has valuation at least
$r+2s+3+2(r-1-s)+2=3r+3$;
its $N^2\sigma_2$ error has valuation at least $4r+2-s\ge3r+4$.
For log degree $j\ge3$, the weighted valuation is at least
$5r-s-1\ge3r+3$, using $jr-v_p(j)\ge3r$ and
$v_p(\sigma_j(pK))\ge\max(0,s-1)$. The exponential remainder
for $A$ has weighted valuation at least $4r+2s+4$.
The second replacement has a harmonic-sum remainder of valuation at least
$2r+s+3+2(r-1-s)=4r+1-s\ge3r+3$.
The higher even-log and exponential terms, after multiplication by
$x^2$, have valuation at least $6r-3-s\ge5r-1\ge3r+3$
by the carry-factor lemma.
All products involving $(A-1)(B-1)$ vanish at the required precision.
Moreover $x^2(E-1)c_0$ has valuation at least $4r+2-s\ge3r+4$.
Thus
\[
 V\equiv\sum_Kx^2O\{(E-1)(a_0+b_0)+b_0+c_0\}
 \pmod{p^{3r+3}}.
\]
The exact alternating identity gives $\sum xO=0$.
Also $\sum_{K<M}O=-1$ and $O_ME_M=\binom{2M}{M}/2$.
Writing $C_M=\binom{2M}{M}$, the terms with $a_0,b_0$ therefore sum to
\[
 \frac{N^2MH_2}{2}\left(\nu(M)+1-\frac{C_M}{2}-3I_1(M)\right).
\]
Here we used $NpM^2=N^2M$ and
$\sum xO(E-1)=I_1(M)$.
The bracket is divisible by $p^3$: this follows from the hypothesis on $\nu(M)$,
$C_M\equiv2\pmod{p^3}$, and
$I_1(M)=\sum xO(E-1)\equiv0\pmod{p^3}$ termwise.
Since $v_p(N^2MH_2)\ge3r$, these terms vanish.

Finally $O_K(M)=(-1)^{K-1}\binom{M-1}{K-1}$ gives
$\sum_{K<M}K O_K(M)=-M$ (the complete first binomial moment is zero),
whereas $\sum_{K<M}O_K(M)/K=0$.
Thus the correction sums exactly to
\[
 \sum_Kx^2O c_0=N^2M^3p^2H_4.
\]
Its valuation is at least $5r-1\ge3r+3$, including $p=5,r=2$.
This completes the proof.
\end{proof}


\section{Completion and consequences}\label{sec:completion}
\begin{proof}[Proof of Theorem~\ref{conj:D}]
The exact alternating identity gives $I_2(n)=nH_{n-1}+R(n)$.
Splitting $R(N)-R(M)$ into terms with $p\nmid k$ and $p\mid k$
gives precisely the unit contribution of Theorem~\ref{thm:unit-all}
and the carry contribution of Lemma~\ref{lem:carry-new}.
The latter lemma applies because Theorem~\ref{thm:nu-all-powers}
gives $\nu(M)\in p^{r+2}\mathbb Z_p\subset p^3\mathbb Z_p$.
Hence, for every $p\ge5,r\ge2$,
\[
R(N)-R(M)\equiv0\pmod{p^{3r+3}}.
\]
The harmonic comparison, including its proved $p=5$ extension,
now gives $\mathcal D_r\equiv0$. Theorem~\ref{thm:towers} gives
the two $B$ rigidity formulas. There is no numerical base case or
unproved induction hypothesis in this argument.

For completeness, the passage to all nonlinear companions is
proved next; the linear statements follow immediately by substitution.
\end{proof}

\begin{lemma}[Polynomial consequences]\label{lem:poly-new}
Put $X(n)=(A(n),A(n-1),B(n),B(n-1))$ and
\[
c=(7/2,-1/2,5/4,-5/4),\qquad X_0=(5,1,3,1).
\]
Let $F$ be a polynomial whose coefficients are $p$-integral.
If $\nabla F(X_0)\mathbin{\cdot}c=0$, then, for $p\ge5,r\ge2$,
\[
F(X(p^r))\equiv F(X(p^{r-1}))\pmod{p^{3r+3}}.
\]
\end{lemma}
\begin{proof}
The preceding theorems give $X(N)-X(M)\equiv\Delta_r c$
modulo $p^{3r+3}$ and $v_p(\Delta_r)\ge3r$.
The exact representations give $X(M)\equiv X_0\pmod{p^3}$:
$J_3,J_4,I_3\in p^3\mathbb Z_p$, $I_2\in p^3\mathbb Z_p$ by
the proof of Theorem~\ref{thm:nu-all-powers}, and
$C_M\equiv2\pmod{p^3}$.
Expand the polynomial at $X(M)$. All terms of degree at least two
in the increment vanish since $6r\ge3r+3$.
Replacing the gradient at $X(M)$ by that at $X_0$ costs a factor
$p^3$ and therefore also vanishes at the desired precision.
The stated gradient identity eliminates the remaining term.
\end{proof}

Applying the lemma to the three polynomials
\[
B^3B^-,\qquad
3^{42}A^{25}-5^{25}B^{42},\qquad
(A^-)^5B^6
\]
proves every nonlinear companion tower. The dot products vanish
exactly, as in the prime-level calculation. No division by $5$ or
by $A(M)$ is made, so the argument includes $p=5$.

\begin{corollary}[Three related OEIS towers]\label{cor:extra-new}
Define
\begin{align*}
u_{12}(n)&=\sum_{k=0}^{n-1}\binom nk^2\binom{n+k}k\binom{n+k-1}k,\\
u_{14}(n)&=\frac14\binom{2n}n^2-\frac{B(n)-3B(n-1)}5,\\
u_{15}(n)&=\sum_{k=0}^{n-1}\binom nk^2\binom{n+k-1}k.
\end{align*}
These are A361712, A361714, and A361715, respectively.
For $r\ge2$ the congruence
$u_i(p^r)\equiv u_i(p^{r-1})\pmod{p^{3r+3}}$
holds for $i=12,15$ and all $p\ge5$, and for $i=14$ and all $p\ge7$.
\end{corollary}
\begin{proof}
The factorizations give $u_{12}(n)=1+J_3(n)+J_4(n)$ and
$u_{15}(n)=1+I_3(n)$ exactly. Their tower congruences have already
been proved. For $u_{14}$, the square-binomial difference divided
by $4$ is congruent to $\Delta_r$, and
$(\alpha_r-3\beta_r)/5\equiv\Delta_r$. Division by $5$ in
this last step is legitimate in the stated range $p\ge7$.
\end{proof}

The definitions and prime ranges in this corollary match the corresponding
OEIS comments \cite{OEIS}. For A361715, the OEIS page credits a proof
of the prime-level statement to Bala; the corollary addresses the tower.

\section{Finite verification}
The accompanying script \texttt{verify\_rigidity\_completion.py} checks exact
binomial identities and modular congruences without floating-point arithmetic.
It passes 673 identity and implementation checks, and 19 tower cases:
$r=2$ for all primes $5\le p\le43$; $p=5$, $2\le r\le5$;
$p=7$, $2\le r\le4$; and $(p,r)=(11,3),(13,3)$.
The largest tested index is $3125$. Its output lists the separate unit, carry,
and rigidity checks. The proof above covers
all primes and all levels in the theorem ranges independently of these tests.
The table below records selected exact valuations from the supplied draft;
the new checks additionally test the completed congruences and moment identities.
\begin{table}[ht]
\centering
\small
\begin{tabular}{@{}lccccccc@{}}
\toprule
$(p,r)$ & $L_r$ & $a(N)-a(M)$ & $\gamma_r-\tfrac72\Delta_r$ & $\delta_r+\tfrac12\Delta_r$ & $\alpha_r-\tfrac54\Delta_r$ & $\beta_r+\tfrac54\Delta_r$ & $\nu(M)$\\
\midrule
$(5,2)$ & 9 & 9 & 9 & 9 & 11 & 12 & 4\\
$(5,3)$ & 12 & 12 & 12 & 12 & 13 & 13 & 5\\
$(5,4)$ & 15 & 15 & 15 & 15 & 16 & 16 & 6\\
$(7,2)$ & 9 & 9 & 9 & 9 & 9 & 9 & 4\\
$(7,3)$ & 12 & 12 & 12 & 12 & 12 & 12 & 5\\
$(11,2)$ & 9 & 9 & 9 & 9 & 9 & 9 & 4\\
$(11,3)$ & 12 & 12 & 12 & 12 & 12 & 12 & 5\\
$(p,2)$, $13\le p\le43$ & 9 & 9 & 9 & 9 & 9 & 9 & 4\\
\bottomrule
\end{tabular}
\caption{$p$-adic valuations of the indicated differences. All five defect columns are covered by the theorems of this revision.
The final column illustrates the lower bound $v_p(\nu(p^s))\ge s+3$.
Exact equality in these samples is not asserted for all primes.}
\label{tab:num}
\end{table}

% Proof fragment to follow apery_rigidity.tex; all dependencies are established there.
\section{A sharp replacement for the enhanced fourth-moment tower}

Put
\[
J_j(n)=\sum_{k=1}^{n-1}(n/k)^jE_k(n)^2,
\qquad u(n)=1+J_4(n),
\qquad d_r=J_4(p^r)-J_4(p^{r-1}).
\]
The factorization in Lemma~\ref{lem:fact} gives
$u(n)=\sum_{k=0}^{n-1}\binom nk^2\binom{n+k-1}k^2$, so this is
OEIS A361713. The original $p^{4r+1}$ proposal and its known counterexample
are recorded in \cite{A361713,Adam}. The proof below uses only the already proved enhanced
$J_3$ stabilization and the elementary carry estimates, not the proposed
stronger congruence for $u$.

\begin{theorem}[Sharp tower law]\label{thm:sharp-J4}
For every prime $p\ge7$ and $r\ge1$,
\[
v_p(d_r)\ge\min(5r,3r+5).
\]
For every $p\ge7$ and $r\ge3$, the first possible defect is given by
\begin{equation}\label{eq:sharp-J4}
d_r\equiv-2p^{3r-1}H_{p-1}^{(2)}J_3(p)
 \pmod{p^{3r+6}}.
\end{equation}
In particular, for the single fixed prime $p=11$,
\[
v_{11}(d_r)=3r+5\qquad\text{for every }r\ge3.
\]
For $p=11$ the first two valuations are $5$ and $10$, respectively.
Thus the exponent $\min(5r,3r+5)$ is attained at every $r\ge1$ by
$p=11$. A proposed congruence $d_r\equiv0\pmod{p^{4r+1}}$ therefore
fails for $p=11$ at \emph{every} $r\ge5$.

For completeness, the small prime $5$ obeys
\[
v_5(d_r)\ge\min(5r-1,3r+4),\qquad
v_5(d_r)=3r+4\quad(r\ge3).
\]
\end{theorem}

We first record the elementary estimates needed at one additional digit
of carry precision. As elsewhere, $H^{(m)}=H_{p-1}^{(m)}$ and
$\sigma_m(X)=\sum_{1\le i<X,\ p\nmid i}i^{-m}$.

\begin{lemma}[Refined quadratic carry]\label{lem:sharp-carry}
Let $p\ge5$, $K\ge1$, and $s=v_p(K)$. Then
\[
\sigma_2(pK)=KH^{(2)}+p^2(K^3-K)H^{(4)}+\mathcal E_K,
\]
where $v_p(\mathcal E_K)\ge s+4$ for $p\ge7$, and
$v_5(\mathcal E_K)\ge s+3$ for $p=5$.
\end{lemma}
\begin{proof}
The convergent carry expansion is
\[
\sigma_2(pK)=\sum_{t\ge0}(-1)^t(t+1)p^tF_t(K)H^{(t+2)},
\qquad F_t(K)=\sum_{a=0}^{K-1}a^t.
\]
Reflection of the reduced residues modulo $p$ gives
\[
H^{(3)}+\tfrac32pH^{(4)}\in p^3\mathbb Z_p.
\]
Indeed, the remaining terms of the reflected expansion start with
$3p^2H^{(5)}$, which lies in $p^3\mathbb Z_p$ because
$H^{(5)}\equiv0\pmod p$ for every $p\ge5$.
Since $v_p(F_1(K))\ge s$, replacing $H^{(3)}$ in the $t=1$ term
introduces an error in $p^{s+4}\mathbb Z_p$. The $t=1$ and $t=2$ terms
then combine exactly as
\[
3p^2\{F_1(K)+F_2(K)\}H^{(4)}
 =p^2(K^3-K)H^{(4)}.
\]
The $t=3$ term has valuation at least $3+2s+1\ge s+4$ because
$F_3(K)=K^2(K-1)^2/4$ and $H^{(5)}\in p\mathbb Z_p$.
For $p\ge7$ the denominator of
\[
F_4(K)=\frac{K(K-1)(2K-1)(3K^2-3K-1)}{30}
\]
is a unit, so its $t=4$ contribution lies in $p^{s+4}\mathbb Z_p$.
For $p=5$ the factor $t+1=5$ compensates the possible denominator
factor $5$, and the same lower bound holds for that contribution.
For $t\ge5$, the carry power-sum bound
$v_p(F_t(K))\ge s-1-v_p(t+1)$, with the factor $t+1$ restored,
gives valuation at least $s+t-1\ge s+4$.
These estimates in fact prove the stated precision for $p\ge7$;
retaining only $s+3$ is sufficient at $p=5$.
\end{proof}



\begin{lemma}[Low odd moment]\label{lem:low-J3}
For $m=p^a$, $a\ge1$, one has
\[
J_3(m)\equiv J_3(p)\pmod{p^9},\qquad
v_p(J_3(m))\ge
\begin{cases}5,&p\ge7,\\4,&p=5.\end{cases}
\]
Also $v_p(J_1(m))\ge1$.
\end{lemma}
\begin{proof}
The identity $a(n)=1-J_3(n)+J_4(n)$, the proved A212334 tower, and
the proved coarse $J_4$ stabilization imply
$J_3(p^b)-J_3(p^{b-1})\in p^{3b+3}\mathbb Z_p$ for $b\ge2$.
This yields the congruence modulo $p^9$.
At the first level,
\[
J_3(p)=p^3H^{(3)}+O(p^5),
\]
because $E_k(p)^2\equiv1\pmod{p^2}$.
For $p\ge7$, $H^{(3)}\in p^2\mathbb Z_p$; for $p=5$,
$H^{(3)}\in p\mathbb Z_p$. These give the asserted valuations.
Finally, every summand of $J_1(m)$ lies in $p\mathbb Z_p$.
\end{proof}

\begin{proof}[Proof of Theorem~\ref{thm:sharp-J4}]
Take $r\ge2$, $N=p^r$, and $M=p^{r-1}$. Split
\[
d_r=U_r+V_r,
\]
where
\[
U_r=N^4\sum_{\substack{k<N\\p\nmid k}}k^{-4}E_k(N)^2
\]
and
\[
V_r=\sum_{K=1}^{M-1}(M/K)^4E_K(M)^2
 (\mathcal B_K^2-1),\qquad
\mathcal B_K=\prod_{\substack{i<pK\\p\nmid i}}(1-N^2/i^2).
\]
These are exact identities.

\emph{The unit terms.}
Always $E_k(N)^2\equiv1\pmod{p^2}$. For $p\ge7$,
$v_p(\sigma_4(N))\ge r$. At $r=2$ this proves
$v_p(U_2)\ge8+2=10$.
At $r=3$ we need one extra digit. With
$q=\lfloor k/p^{r-1}\rfloor$ and $h_q=\sum_{a=1}^qa^{-2}$, the
already established top-level expansion is
\[
E_k(N)^2\equiv1-2p^2h_q\pmod{p^4}.
\]
Each complete top block has
\[
\sum_{\substack{qp^{r-1}<k<(q+1)p^{r-1}\\p\nmid k}}k^{-4}
\equiv p^{r-2}H^{(4)}\equiv0\pmod p.
\]
Consequently the weighted unit sum is congruent to $\sigma_4(N)$
modulo $p^3$. For $r=3$ this yields $v_p(U_3)\ge15$.
For $r=4$, $E_k(N)^2\equiv1\pmod{p^2}$ gives $v_p(U_4)\ge18$;
for $r=5$, reduction modulo $p$ gives $v_p(U_5)\ge21$;
and for $r\ge6$, the factor $N^4$ alone gives
$4r\ge3r+6$. We have proved
\begin{equation}\label{eq:unit-sharp}
v_p(U_2)\ge10,\qquad v_p(U_r)\ge3r+6\quad(r\ge3,\ p\ge7).
\end{equation}

For $p=5$, one has the general bound
$v_p(\sigma_m(p^r))\ge r-1$. To see this, lift each unit residue
modulo $p^{r-1}$ to its $p$ representatives modulo $p^r$ and reduce
modulo $p^{r-1}$, then induct on $r$.
Thus at $r=2$, $v_5(U_2)\ge9$. At $r=3$, use
$E_k(N)^2\equiv1\pmod{p^2}$ and $\sigma_4(5^3)\in5^2\mathbb Z_5$
to get $v_5(U_3)\ge14$. At $r=4$, reduction modulo $5$ gives
$v_5(U_4)\ge17$. For $r\ge5$, $4r\ge3r+5$. Hence
\begin{equation}\label{eq:unit-sharp-five}
v_5(U_2)\ge9,\qquad v_5(U_r)\ge3r+5\quad(r\ge3).
\end{equation}

\emph{Carry errors and their leading term.}
Put $s=v_p(K)\le r-2$. The proved carry-factor estimate gives
\[
\mathcal B_K^2-1=-2N^2\sigma_2(pK)+O(p^{4r+s-1}).
\]
After multiplication by $(M/K)^4E_K(M)^2$, the discarded error
has valuation at least
\[
4(r-1-s)+4r+s-1=8r-5-3s\ge5r+1.
\]
For $p\ge7$, insert Lemma~\ref{lem:sharp-carry}.
Its remainder after multiplication has valuation at least
\[
2r+4(r-1-s)+s+4=6r-3s\ge3r+6.
\]
Using the exact identities
\[
\sum_K(M/K)^4K E_K(M)^2=MJ_3(M),\qquad
\sum_K(M/K)^4K^3E_K(M)^2=M^3J_1(M),
\]
we obtain for $r\ge3$
\begin{align*}
V_r\equiv{}&-2N^2M H^{(2)}J_3(M)\\
&-2N^2p^2H^{(4)}\{M^3J_1(M)-MJ_3(M)\}
\pmod{p^{3r+6}}.
\end{align*}
Here $5r+1\ge3r+6$, so the carry-factor error also vanishes.
The $J_1$ term has valuation at least $5r+1$ and the second $J_3$
term at least $3r+7$, using Lemma~\ref{lem:low-J3} and
$H^{(4)}\in p\mathbb Z_p$. They therefore vanish. Replacing
$J_3(M)$ with $J_3(p)$ is legitimate modulo $p^{3r+6}$ by the
same lemma. Together with \eqref{eq:unit-sharp}, this proves
\eqref{eq:sharp-J4}, and hence $v_p(d_r)\ge3r+5$ for $r\ge3$.

For $r=2$, the weaker carry approximation
$\sigma_2(pK)=KH^{(2)}+O(p^{s+3})$ suffices: it gives
\[
V_2=-2N^2M H^{(2)}J_3(M)+O(p^{3r+5})\in p^{11}\mathbb Z_p.
\]
With $v_p(U_2)\ge10$ this proves the assertion at $r=2$.
At $r=1$, $d_1=J_4(p)=p^4H^{(4)}+O(p^6)\in p^5\mathbb Z_p$.
This proves the full lower bound for $p\ge7$.

At $p=5$, use the $s+3$ precision in
Lemma~\ref{lem:sharp-carry}. Its weighted error has valuation at least
$6r-1-3s\ge3r+5$. The two correction terms above have valuations
at least $5r$ and $3r+5$, respectively, because now
$v_5(H^{(4)})\ge0$ and $v_5(J_3(M))\ge4$.
Thus, for $r\ge3$, the same argument with
\eqref{eq:unit-sharp-five} gives
\begin{equation}\label{eq:sharp-five}
d_r\equiv-2\cdot5^{3r-1}H_4^{(2)}J_3(5)
 \pmod{5^{3r+5}}.
\end{equation}
At $r=2$, the weaker carry approximation with error $s+2$ gives
$v_5(V_2)\ge3r+4=10$, while $v_5(U_2)\ge9$.
At $r=1$, integrality of $H^{(4)}$ gives $v_5(d_1)\ge4$.
These prove the stated $5$-adic lower bounds.

\emph{Exact sharpness certificates.}
For $p=11$, direct finite rational arithmetic gives
\[
H_{10}^{(2)}=\frac{1968329}{1270080},\qquad
J_3(11)=97380253269607=11^5\cdot604654757.
\]
Consequently
\[
H_{10}^{(2)}/11\equiv10\pmod{11},\qquad
J_3(11)/11^5\equiv3\pmod{11}.
\]
Dividing \eqref{eq:sharp-J4} by $11^{3r+5}$ therefore yields
\[
d_r/11^{3r+5}\equiv-2\cdot10\cdot3\equiv6\pmod{11},
\]
which is nonzero for every $r\ge3$ and proves exact sharpness.
To verify sharpness at the first two levels as well, note that
\[
d_1\equiv p^4H^{(4)}\pmod{p^6},\qquad
 d_2\equiv p^9H^{(4)}\pmod{p^{11}}\quad(p\ge7).
\]
For the second congruence, the top-block argument above shows
$U_2\equiv p^8\sigma_4(p^2)\pmod{p^{11}}$, while $V_2\in p^{11}\mathbb Z_p$.
Expansion in the $p$ complete blocks gives
$\sigma_4(p^2)\equiv pH^{(4)}\pmod{p^3}$: its linear correction
has valuation at least $1+v_p(F_1(p))+v_p(H^{(5)})\ge3$,
its quadratic correction has valuation at least
$2+v_p(F_2(p))\ge3$, and all later terms have at least this valuation
by the same power-sum bounds.
For $p=11$, the exact finite value
\[
H_{10}^{(4)}=\frac{43635917056897}{40327580160000}
\equiv44\pmod{11^2}
\]
therefore gives $d_1/11^5\equiv d_2/11^{10}\equiv4\pmod{11}$.
For $p=5$,
\[
H_4^{(2)}=\frac{205}{144},\qquad J_3(5)=180625=5^4\cdot289.
\]
The corresponding residues are both $4$ modulo $5$, so
\eqref{eq:sharp-five} gives
\[
d_r/5^{3r+4}\equiv-2\cdot4\cdot4\equiv3\pmod5
\qquad(r\ge3).
\]
This proves the claimed exact $5$-adic law.
\end{proof}

% Self-contained proof fragment for the smooth-index claim on A357567.
\section{An infinite counterexample family for the smooth-index claim}

Retain the definitions
\[
 A(n)=\sum_{k=0}^n\binom nk^2\binom{n+k}k^2,
 \qquad
 B(n)=\sum_{k=0}^n\binom nk^2\binom{n+k}k,
\]
and put
\[
 F(n)=3^{42}A(n)^{25}-5^{25}B(n)^{42}.
\]
The fourth conjecture recorded on OEIS A357567 \cite{A357567} asserts
$F(n)\equiv0\pmod{n^5}$ for $n=3^i5^j$ with $i,j\ge0$.
This statement is false. In fact, the obstruction persists along an entire
infinite family, and its proof only needs a binomial congruence modulo $25$.

\begin{lemma}\label{lem:transfer-babbage-five}
For integers $a\ge b\ge0$,
\[
 \binom{5a}{5b}\equiv\binom ab\pmod{25}.
\]
Consequently, for every integer $m\ge0$,
\[
 A(5m)\equiv A(m)\pmod{25},\qquad
 B(5m)\equiv B(m)\pmod{25}.
\]
\end{lemma}
\begin{proof}
Let
\[
 U(a)=\prod_{\substack{1\le t\le5a\\5\nmid t}}t,
 \qquad U(0)=1.
\]
For every $j\ge0$, expansion in $\mathbb Z_5$ gives
\[
 \prod_{u=1}^4(5j+u)
 \equiv24\left(1+5j\sum_{u=1}^4\frac1u\right)
 \equiv24\pmod{25},
\]
because $\sum_{u=1}^4u^{-1}=25/12\equiv0\pmod5$.
Hence $U(a)\equiv24^a\pmod{25}$. Splitting the factorials into multiples
and nonmultiples of $5$ gives the exact equality
\[
 \binom{5a}{5b}
 =\binom ab\frac{U(a)}{U(b)U(a-b)}.
\]
The quotient on the right is a $5$-adic unit congruent to $1$ modulo $25$,
which proves the binomial congruence, without dividing by $\binom ab$.

In either Ap\'ery sum at $5m$, the terms with $5\nmid k$ vanish modulo
$25$, since
\[
 \binom{5m}{k}=\frac{5m}{k}\binom{5m-1}{k-1}\in5\mathbb Z_5
\]
and this factor is squared. For $k=5l$, apply the proved congruence to
$\binom{5m}{5l}$ and $\binom{5(m+l)}{5l}$. Their contributions reproduce
the corresponding sums at $m$ modulo $25$.
\end{proof}

\begin{proposition}[A residue formula along powers of five]
\label{prop:transfer-residue-five}
Suppose $m\ge1$ and $5\mid A(m)$. For every integer $r\ge0$, the integer
$F(5^rm)$ is divisible by $5^{25}$ and satisfies
\[
 \frac{F(5^rm)}{5^{25}}
 \equiv 4\frac{A(m)}5-B(m)^2\pmod5.
\]
If the right side is nonzero modulo $5$, then
$v_5(F(5^rm))=25$ for every $r\ge0$.
\end{proposition}
\begin{proof}
Iterating Lemma~\ref{lem:transfer-babbage-five} gives
$A(5^rm)\equiv A(m)\pmod{25}$ and
$B(5^rm)\equiv B(m)\pmod5$. Therefore
\[
 \frac{F(5^rm)}{5^{25}}
 =3^{42}\left(\frac{A(5^rm)}5\right)^{25}-B(5^rm)^{42}.
\]
Modulo $5$, the polynomial identities $x^{25}=x$ and $x^{42}=x^2$
hold for every residue $x$, including $x=0$, and $3^{42}\equiv4$.
Substitution gives the assertion.
\end{proof}

\begin{theorem}[Infinite counterexamples to A357567, conjecture 4]
\label{thm:transfer-counterexample-five}
For every integer $r\ge0$,
\[
 \boxed{v_5\!\left(F(3\cdot5^r)\right)=25.}
\]
In particular, every $n=3\cdot5^r$ with $r\ge6$ is a counterexample to
$F(n)\equiv0\pmod{n^5}$. One explicit counterexample is
$n=3\cdot5^6=46875$.
\end{theorem}
\begin{proof}
The defining sums give $A(3)=1445$ and $B(3)=147$. Thus
\[
 \frac{A(3)}5\equiv4\pmod5,\qquad B(3)\equiv2\pmod5,
\]
and Proposition~\ref{prop:transfer-residue-five} yields
\[
 \frac{F(3\cdot5^r)}{5^{25}}
 \equiv4\cdot4-2^2\equiv2\pmod5.
\]
This proves the exact valuation. For $r\ge6$,
$v_5((3\cdot5^r)^5)=5r>25$, so divisibility by $n^5$ is impossible.
In particular, the explicit certificate for $n=46875$ is
\[
 F(46875)\equiv2\cdot5^{25}
 =596046447753906250\not\equiv0\pmod{5^{26}}.
\]
No evaluation of the full enormous integer $F(46875)$ is needed.
\end{proof}

\begin{remark}[Scope]
This refutes the smooth-index assertion, including its version with both
$i$ and $j$ positive. It does not contradict the adjacent prime-power
congruences for $F(p^r)-F(p^{r-1})$: those statements compare successive
values and do not require $F(p^r)$ itself to tend to zero $p$-adically.
The proof above does not assert that $46875$ is the least counterexample,
and it is independent of the enhanced-rigidity theorem.
\end{remark}

\section{Bala's second conjecture for A361717}
The official OEIS source file
\cite{A361717}
(revision 46, 2026-09-08) states Conjecture~2 as
\[
 \nu(p-1)\equiv1-2p-p^2\pmod{p^3}
 \qquad\text{for every prime }p\ne3.
\]
It is a congruence at $p-1$, not a prime-power tower.

\begin{lemma}[An adjacent exact identity]
For every integer $n\geq2$,
\[
 (n-2)\nu(n-1)=(5n-2)B(n-1)-3n\nu(n).
\]
\end{lemma}
\begin{proof}
Let $u_k=\binom{n-1}{k}^2\binom{n+k}{k}$ for $0\leq k\leq n-1$
and set $u_n=0$. The summands of $\nu(n-1)$ and $B(n-1)$ are
respectively
\[
 \frac{n(n-1-k)^2}{(n-1)^2(n+k)}u_k,
 \qquad \frac{n}{n+k}u_k.
\]
The rational certificate
\[
 G(k)=-\frac{nk^3}{(n-1)^2(n+k)}
\]
satisfies
\begin{align*}
 G(k+1)u_{k+1}-G(k)u_k
 =\left(\frac{n(n-2)(n-1-k)^2}{(n-1)^2(n+k)}
       -\frac{n(5n-2)}{n+k}+3n\right)u_k.
\end{align*}
This follows by substituting
$u_{k+1}/u_k=(n-1-k)^2(n+k+1)/(k+1)^3$ and clearing denominators.
Summing telescopes, and the boundary terms vanish.
\end{proof}

\begin{theorem}[Bala's A361717 Conjecture 2]
For every prime $p\ne3$,
\[
 \nu(p-1)\equiv1-2p-p^2\pmod{p^3}.
\]
\end{theorem}
\begin{proof}
Suppose first that $p\geq7$. The prime-level rigidity formulas imply
$B(p-1)\equiv1\pmod{p^3}$, and the zeroth-moment theorem gives
$\nu(p)\equiv0\pmod{p^4}$. The adjacent exact identity therefore gives
\[
 \nu(p-1)\equiv\frac{5p-2}{p-2}\pmod{p^3}.
\]
The exact rational equality
\[
 \frac{5p-2}{p-2}-(1-2p-p^2)=\frac{p^3}{p-2}
\]
finishes the proof. At $p=5$, one has $\nu(4)=216$ and
$216-(1-10-25)=250$, divisible by $5^3$. At $p=2$, one has
$\nu(1)=1$ and $1-(1-4-4)=8$. At the excluded prime $p=3$,
$\nu(2)-(1-6-9)=18$, which is not divisible by $27$.
\end{proof}

\begin{remark}[Scope]
This is a short consequence of the adjacent identity and the proved
prime-level estimates. The source entry states precisely this prime range.
Novelty relative to unpublished work has not been established.
\end{remark}

\begin{corollary}[A fifth-order refinement]
For every prime $p\geq7$, writing
$\Delta_1=\binom{2p}{p}-2$, one has
\[
 \nu(p-1)\equiv\frac{5p-2}{p-2}
       \left(1-\frac54\Delta_1\right)\pmod{p^5}.
\]
\end{corollary}
\begin{proof}
Use $B(p-1)\equiv1-5\Delta_1/4\pmod{p^5}$ in the adjacent
identity and note that $p\nu(p)\equiv0\pmod{p^5}$.
\end{proof}

\section{A normalized zeroth-moment tower}
By the completed $I_2$ and $I_3$ tower theorems proved above, including $p=5$, for every prime $p\geq5$ and $r\geq2$, with
$N=p^r$, $M=p^{r-1}$, and
$\Delta_r=\binom{2N}{N}-\binom{2M}{M}$,
\[
 (N-1)\nu(N)-p(M-1)\nu(M)
       \equiv N\Delta_r\pmod{p^{4r+3}}.
\]
Indeed the exact identities imply
\[
 \frac{n-1}{n}\nu(n)=\frac12\binom{2n}{n}-1+2I_2(n)-I_3(n).
\]
Taking adjacent-level differences, substituting
$\Delta I_2\equiv\Delta_r/4$ and $\Delta I_3\equiv0$
modulo $p^{3r+3}$, and multiplying by $N$ proves the claim.

\section{The higher target: the remaining Swisher boundary}
Swisher's $(C.3)$ conjecture has the form
\[
S_r=\sum_{k=0}^{(p^r-1)/2}(4k+1)\frac{(1/2)_k^4}{k!^4},\qquad
S_r\equiv pS_{r-1}\pmod{p^{4r}}\quad(p>3).
\]
Guo proves this family modulo $p^{3r}$ \cite{Guo}; that source does not
give the full displayed modulus. The present transfer does not complete
the $p^{4r}$ assertion.

Here is the exact obstruction to simply importing a homogeneous moment
tower. Put $w_k=(1/2)_k^4/k!^4$ and, for odd $n\ge1$, put
\[
T(n)=\sum_{k=0}^{(n-1)/2}(4k+1)k(k+1/2)w_k,\qquad
D(n)=(n+1)^4w_{(n+1)/2}.
\]
The recurrence $16(k+1)^4w_{k+1}=(2k+1)^4w_k$ implies
\[
16(k+1)^4w_{k+1}-16k^4w_k
=(4k+1)(8k^2+4k+1)w_k.
\]
Summing gives $S(n)+8T(n)=D(n)$ exactly, with $S(n)$ defined by the
same truncated sum as $S_r$ at $n=p^r$. Moreover,
\[
D(n)=\frac{n^4}{256^{(n-1)/2}}
       \binom{n-1}{(n-1)/2}^{\!4}.
\]
For $n=p^r$, the binomial coefficient is a $p$-adic unit: its residue
is $(-1)^{(p^r-1)/2}$, by the coefficient identity
$(1+x)^{p^r-1}\equiv\sum_{k=0}^{p^r-1}(-x)^k\pmod p$.
Thus $v_p(D(p^r))=4r$.
Consequently $(C.3)$ is equivalent to the inhomogeneous requirement
\[
8\bigl(T(p^r)-pT(p^{r-1})\bigr)
\equiv-pD(p^{r-1})\pmod{p^{4r}}.
\]
The source term on the right has valuation exactly $4r-3$ and cannot
be discarded. The proof of this last congruence remains outstanding.

\section{Final scope and verification}\label{sec:status}
The original rigidity packet is complete in its stated prime ranges.
The two separate assertions left unresolved in the preceding draft now
have definite outcomes: A357567 conjecture~4 is false, with an infinite
counterexample family; A361717 conjecture~2 is true for its entire stated
range. The sharp A361713 theorem gives a stronger uniform law whose
exponent is optimal at every level.

The scripts \texttt{verify\_transfer\_sharp.py} and
\texttt{verify\_smooth\_counterexamples.py}, with their exact outputs,
are included beside the prior checker. They use integer modular arithmetic
and rational arithmetic, not floating point. The scripts specify the finite
ranges checked; every infinite-family assertion has a written proof above.
The Swisher $(C.3)$ result is a reduction only, not a claimed solution.
The source review checked the stated OEIS entries, the public campaign,
and the cited primary papers. It does not establish priority, and the
proof draft has not undergone external peer review or formal verification.

\begin{thebibliography}{99}
\bibitem{Ap} R.~Ap\'ery, Irrationalit\'e de $\zeta(2)$ et $\zeta(3)$, \emph{Ast\'erisque} 61 (1979), 11--13.
\bibitem{Be} F.~Beukers, Some congruences for the Ap\'ery numbers, \emph{J. Number Theory} 21 (1985), 141--155.
\bibitem{Co} M.~J.~Coster, \emph{Supercongruences}, Ph.D. thesis, Universiteit Leiden, 1988.
\bibitem{Ge} I.~Gessel, Some congruences for Ap\'ery numbers, \emph{J. Number Theory} 14 (1982), 362--368.
\bibitem{Gr} A.~Granville, Arithmetic properties of binomial coefficients I: binomial coefficients modulo prime powers, \emph{CMS Conf. Proc.} 20 (1997), 253--275.
\bibitem{St} A.~Straub, Multivariate Ap\'ery numbers and supercongruences of rational functions, \emph{Algebra Number Theory} 8 (2014), 1985--2008.
\bibitem{Zh} J.~Zhao, Wolstenholme type theorem for multiple harmonic sums, \emph{Int. J. Number Theory} 4 (2008), 73--106.
\bibitem{OEIS} The On-Line Encyclopedia of Integer Sequences, entries A005258, A005259, A212334, A352655, A357506, A357567, A357956, A357959, A361712, A361714, A361715, A361717 (comments by P.~Bala, 2021--2023), \url{https://oeis.org}.
\bibitem{FC} Google DeepMind, \emph{Formal Conjectures}, file \texttt{FormalConjectures/OEIS/212334.lean}, snapshot of 2026-11-11, \url{https://github.com/google-deepmind/formal-conjectures}.
\bibitem{Cam} GitHub user \texttt{rbajaj5}, \emph{a183068-supercongruence} (public repository, 2026), notes ``The enhanced prime congruence for A212334'' and ``The Ap\'ery enhanced-congruence packet'', \url{https://github.com/rbajaj5/a183068-supercongruence}.

\bibitem{A361713} P.~Bala and contributors, OEIS A361713, including the
2026 correction to Conjecture 2, \url{https://oeis.org/A361713}.
\bibitem{A357567} P.~Bala and contributors, OEIS A357567, conjectures 1--4,
\url{https://oeis.org/A357567}.
\bibitem{A361717} P.~Bala and contributors, OEIS A361717, official source
revision 46 (September 8, 2026),
\url{https://github.com/oeis/oeisdata/blob/main/seq/A361/A361717.seq}.
\bibitem{Adam} T.~Adamczewski, \emph{OEIS Open: How many conjectures can
language models turn into theorems?}, arXiv:2608.11941v2 (2026),
\url{https://arxiv.org/abs/2608.11941}.
\bibitem{Guo} V.~J.~W.~Guo, $q$-Analogues of Dwork-type supercongruences,
\emph{J. Math. Anal. Appl.} 487 (2020), 124022,
\url{https://arxiv.org/abs/1910.07551}.

\end{thebibliography}

\end{document}
